\chapter{Adaptive Memory Routing and Memory Management}\label{ch:routing}

The hybrids of \cref{ch:hybrid} combine two kinds of working memory: a recurrent state of fixed size
(\cref{ch:linear,ch:delta}) and a softmax-attention KV cache that grows with every token
(\cref{ch:foundations}). They make the split once, at design time. Some layers or heads are attention
and the rest are recurrent, often in a ratio such as $3{:}1$. Two questions are left open. First,
\emph{which tokens} need the expensive, high-resolution memory? A fixed layout sends every token
through the attention layers and writes every token into their caches, although most tokens are
predicted well by the cheap recurrent path. Second, once something is in working memory, \emph{how
should that memory be managed}? A cache entry can be stored in fewer bits, a recurrent state can be
copied and kept for later, and a stored entry can be rewritten after the model has seen more.

This chapter covers both questions and has two halves. \textbf{Part A} (adaptive routing) covers
systems that decide per token which memory pathway is used. AMOR~\citep{zheng2026amor} runs a
recurrent language model and turns on an attention refinement only where the model's own predictive
entropy is high. HAM~\citep{lufkin2026ham} lets a recurrent state see every token, but writes into a
KV cache only the tokens the recurrent state predicts poorly. The survey~\citep{zhoubian2026memory}
calls this signal-gated routing an emerging direction in hybrid design: the question is not only
\emph{where} attention sits, but \emph{which tokens} receive high-resolution storage or refinement.
\textbf{Part B} (memory management) treats the working memory itself as something to engineer.
CommVQ~\citep{li2025commvq} vector-quantizes the KV cache to one or two bits per channel, using a
codebook that commutes with the rotary position embedding. Memory Caching~\citep{behrouz2026memorycaching}
keeps snapshots of a recurrent memory at segment boundaries, so a recurrent model's memory grows
with the sequence. Bottlenecked Transformers~\citep{oomerjee2026bottlenecked} rewrite KV entries in
place at reasoning-step boundaries, a form of consolidation.

\paragraph{Roadmap.} \Cref{sec:c13-signals} builds the background for Part A: predictive entropy,
surprise and prediction error, admission and eviction policies, and how to train through hard
decisions. \Cref{sec:c13-amor,sec:c13-ham} cover AMOR and HAM. \Cref{sec:c13-mgmt-bg} builds the
background for Part B: the size of the KV cache, scalar and vector quantization (codebooks, $k$-means,
EM, residual and additive VQ), a recap of RoPE and why quantization must respect it, and
checkpointing a recurrent state. \Cref{sec:c13-commvq,sec:c13-mc,sec:c13-bottleneck} cover CommVQ,
Memory Caching and Bottlenecked Transformers. \Cref{sec:c13-compare} compares the five systems.

% ===========================================================================
\section{Background I: signals and policies for adaptive memory}\label{sec:c13-signals}
% ===========================================================================

\subsection{What a hybrid layer costs}\label{sec:c13-cost}

Recall the two working memories. A softmax-attention head keeps every past key and value. At step
$t$ its cache holds $t$ keys $\vk_s\in\R^{d_k}$ and $t$ values $\vv_s\in\R^{d_v}$, and producing one
output costs $\bigO(t(d_k+d_v))$ multiply-adds. A linear-attention or delta-rule head keeps a matrix
$\mS_t\in\R^{d_v\times d_k}$, so its memory and per-token cost are $\bigO(d_kd_v)$, independent of $t$.
Consider a model with $L$ layers, of which a fraction $\rho$ are attention layers with $H_{kv}$ key-value
heads each, and the rest are recurrent layers with $H$ heads. With $b$ bytes per stored number, the
memory after $T$ tokens is
\begin{equation}\label{eq:c13-hybrid-mem}
  \underbrace{\rho L \cdot 2\,T\,H_{kv}\,d_k\cdot b}_{\text{KV caches, grows with }T}
  \;+\; \underbrace{(1-\rho)L\cdot H\,d_k d_v\cdot b}_{\text{recurrent states, constant}},
\end{equation}
where we took $d_v=d_k$ in the KV term. A fixed-ratio hybrid (\cref{ch:hybrid}) lowers $\rho$, but
every attention layer still pays for all $T$ tokens. An \emph{adaptive} design replaces $T$ in the
first term by the number of tokens actually admitted, $N_T=\sum_{t\le T} g_t$, where $g_t\in\{0,1\}$ is
a per-token decision. A design that saves compute but not memory instead replaces the per-token
attention cost $\bigO(t\,d_k)$ by $g_t\cdot\bigO(t\,d_k)$. AMOR is of the second kind and HAM of the
first, and the decision $g_t$ is the subject of this section.

\subsection{Predictive entropy}\label{sec:c13-entropy}

A language model ends each position with logits $\vz_t\in\R^{\abs{\mathcal{V}}}$ and the next-token
distribution $\vec{p}_t=\softmax(\vz_t)$. The \emph{entropy} of this distribution~\citep{c13-shannon1948} is
\begin{equation}\label{eq:c13-entropy}
  \mathrm{H}(\vec{p}_t) = -\sum_{i=1}^{\abs{\mathcal{V}}} p_{t,i}\log p_{t,i}
  \qquad(\text{with } 0\log 0 \defeq 0).
\end{equation}
It is the \emph{expected surprise}: if the next token $x_{t+1}$ is drawn from $\vec{p}_t$, then
$\E[-\log p_{t,x_{t+1}}]=\mathrm{H}(\vec{p}_t)$. It measures how unsure the model is \emph{before} it sees
the token, so a model can use it at inference time to decide how much work to spend on position~$t$.

\paragraph{Bounds and normalization.} Each term $-p\log p$ is non-negative for $p\in[0,1]$, so
$\mathrm{H}\ge 0$, with equality only for a one-hot distribution. For the upper bound, use Jensen's inequality on the
concave function $\log$:
\[
  \mathrm{H}(\vec{p}) = \sum_i p_i \log\frac{1}{p_i} \le \log\Bigl(\sum_i p_i\cdot\frac{1}{p_i}\Bigr)
  = \log \abs{\{i: p_i>0\}} \le \log\abs{\mathcal{V}},
\]
with equality for the uniform distribution. Dividing by $\log\abs{\mathcal{V}}$ therefore gives a
\emph{normalized entropy} $\bar{\mathrm{H}}_t = \mathrm{H}(\vec{p}_t)/\log\abs{\mathcal{V}}\in[0,1]$, which can be compared
across vocabularies and is what AMOR thresholds.

\begin{example}[A four-token vocabulary]\label{ex:c13-entropy}
Let $\abs{\mathcal{V}}=4$ and $\vec{p}=(0.7,0.1,0.1,0.1)$. Then
$\mathrm{H}=-0.7\log 0.7-3\cdot 0.1\log 0.1 = 0.2497+0.6908=0.9404$ nats, and $\log 4=1.3863$, so
$\bar{\mathrm{H}}=0.678$. For the two-way tie $\vec{p}=(0.4,0.4,0.1,0.1)$ we get $\mathrm{H}=1.1935$ and
$\bar{\mathrm{H}}=0.861$. The tie is ``more uncertain'' even though its top probability is only a little lower.
Entropy looks at the whole distribution, not only at the top probability.
\end{example}

\paragraph{Computing it stably.} Write $\log p_i = z_i-\operatorname{lse}(\vz)$ with
$\operatorname{lse}(\vz)=\log\sum_j e^{z_j}$. Substituting into \eqref{eq:c13-entropy},
\begin{equation}\label{eq:c13-entropy-lse}
  \mathrm{H}(\vec{p}) = -\sum_i p_i\bigl(z_i-\operatorname{lse}(\vz)\bigr) = \operatorname{lse}(\vz)-\sum_i p_i z_i,
\end{equation}
because $\sum_i p_i=1$. The log-sum-exp is computed with the usual max-shift, and no $\log 0$ appears.
In practice the logits are cast to fp32 first, because $\abs{\mathcal{V}}$ is large and low-precision
sums lose accuracy.

\paragraph{What it costs.} The entropy needs the full logit vector, i.e.\ a product with the
$\abs{\mathcal{V}}\times d$ output matrix. For $d=768$ and $\abs{\mathcal{V}}=128256$ (the AMOR 180M
configuration, \cref{sec:c13-amor}) this is $768\cdot 128256\approx 9.85\times 10^{7}$ multiply-adds per
token. One attention read over $t$ cached tokens costs about $2td$ multiply-adds, so the gate costs more
than the attention it may skip unless the context is longer than about $64{,}000$ tokens. AMOR's answer, a small
fitted network that predicts the entropy, appears in \cref{sec:c13-amor-router}.

\paragraph{Uncertainty is not the same as need.} High entropy has two sources. The model may lack
information that is present in the context (it ``forgot'' a name seen 2000 tokens ago). This
uncertainty can be reduced by looking back. Or the next token may be genuinely unpredictable (the
first word of a new sentence). This uncertainty cannot be reduced by any amount of memory. The
entropy signal cannot tell the two apart. Systems that gate on it accept some wasted refinement on the
second kind.

\subsection{Surprise and prediction error}\label{sec:c13-surprise}

Entropy is a property of a prediction. \emph{Surprise} is a property of an outcome: once $x_{t+1}$ is
known, its surprise is $-\log p_{t,x_{t+1}}$. The same idea applies inside a memory. An associative
memory $\mS_{t-1}$ ``predicts'' the value of a new key as $\hat{\vv}_t=\mS_{t-1}\vk_t$. The
\emph{prediction error}, or innovation,
\begin{equation}\label{eq:c13-innovation}
  \ve_t = \vv_t-\mS_{t-1}\vk_t ,
\end{equation}
is the part of the new association the memory does not already contain. This is exactly the quantity
the delta rule writes (\cref{ch:delta}): with write strength $\beta_t$,
$\mS_t=\mS_{t-1}+\beta_t\ve_t\vk_t\T$. If we model the value as Gaussian around the memory's
prediction, $\vv_t\sim\mathcal{N}(\mS_{t-1}\vk_t,\sigma^2\mI)$, then
\[
  -\log p(\vv_t\mid\vk_t) = \frac{\norm{\ve_t}^2}{2\sigma^2} + \frac{d_v}{2}\log(2\pi\sigma^2),
\]
so the squared prediction error \emph{is} the surprise up to an affine map. A scale-free version that
does not depend on the size of $\vv_t$ is the normalized error
\begin{equation}\label{eq:c13-normerr}
  \varepsilon_t = \frac{\norm{\vv_t-\mS_{t-1}\vk_t}^2}{\norm{\vv_t}^2+\epsilon}.
\end{equation}
$\varepsilon_t=0$ means the memory already recalls $\vv_t$ perfectly, and $\varepsilon_t\approx 1$ means it
recalls nothing useful (for example, $\mS_{t-1}\vk_t\approx\bm{0}$). Titans (\cref{ch:ttt}) uses the
gradient of such an error to modulate parametric writes. HAM (\cref{sec:c13-ham}) uses it to decide
which tokens deserve an exact copy in a KV cache.

\subsection{Admission and eviction policies}\label{sec:c13-policies}

Let the store after step $t$ be a set of positions $\mathcal{C}_t\subseteq\{1,\dots,t\}$ whose
keys and values are kept. A \emph{policy} decides how $\mathcal{C}_t$ evolves. Two kinds of decisions
are useful to separate.
\begin{itemize}
  \item An \textbf{admission} decision is made once, when token $t$ is written:
        $\mathcal{C}_t=\mathcal{C}_{t-1}\cup\{t\}$ if $g_t=1$ and $\mathcal{C}_t=\mathcal{C}_{t-1}$ otherwise.
        A token that is not admitted can never be recalled exactly later, although a summary
        of it may survive in a recurrent state.
  \item An \textbf{eviction} decision removes an entry already in the store, usually because a
        budget $\abs{\mathcal{C}_t}\le B$ is exceeded.
\end{itemize}
Classical cache policies are eviction policies: FIFO (a sliding window), least-recently-used, and the
offline optimum of \citet{c13-belady1966}, which evicts the entry whose next use is furthest in the
future and is therefore not implementable online. LLM-specific policies keep a few initial ``attention
sink'' tokens plus a recent window~\citep{xiao2023streamingllm}, or keep ``heavy hitters'' with large
accumulated attention~\citep{c13-zhang2023h2o}. The survey (Table~II) lists admission, eviction and
consolidation as one update-rule family, with ``irreversible loss and retrieval bias'' as its main
risk~\citep{zhoubian2026memory}.

\paragraph{Threshold policies and budget control.} The simplest admission policy compares a score
$s_t$ with a threshold: $g_t=\mathbf{1}[s_t>\tau]$. The expected number of admitted tokens is
\begin{equation}\label{eq:c13-expected-cache}
  \E\bigl[\abs{\mathcal{C}_T}\bigr] = \sum_{t=1}^{T}\Pr(s_t>\tau),
\end{equation}
so $\tau$ directly controls the growth rate of the cache: if the scores are identically distributed,
the cache grows like $rT$ with \emph{admission rate} $r=\Pr(s>\tau)$. To hit a target rate $r$, choose
$\tau$ as the $(1-r)$-quantile of the score distribution. Quantiles drift during training, so systems
track them with an exponential moving average (EMA). AMOR uses the median plus a multiple of the
standard deviation. If the scores were Gaussian, $s\sim\mathcal{N}(\mu,\sigma^2)$, then the median is
$\mu$ and $\tau=\mu+\kappa\sigma$ gives
\begin{equation}\label{eq:c13-fire-rate}
  r = \Pr\Bigl(\frac{s-\mu}{\sigma}>\kappa\Bigr) = 1-\Phi(\kappa),
\end{equation}
with $\Phi$ the standard normal CDF. For $\kappa=0.2$ this is $r\approx 0.42$. A small NumPy
simulation of the EMA tracker (in our scratch space) settles at the same value. Real entropy
distributions are skewed, so this is a guide, not a guarantee.

\subsection{Training through hard decisions}\label{sec:c13-hard}

The indicator $\mathbf{1}[s>\tau]$ has zero derivative almost everywhere, so a gate cannot be learned
by plain backpropagation. Four standard responses appear in this chapter.
\begin{enumerate}
  \item \textbf{Detach the gate.} Treat $g_t$ as a constant and train only the computation it
        selects. Gradients reach the model through the outputs it gates, not through the decision.
        AMOR does this.
  \item \textbf{Straight-through estimator}~\citep{c10-bengio2013ste}. Use the hard value in the
        forward pass and pretend it was the identity (or a sigmoid) in the backward pass:
        $y = y_{\text{hard}} + y_{\text{soft}} - \operatorname{sg}(y_{\text{soft}})$, where $\operatorname{sg}$
        stops the gradient. The forward value is $y_{\text{hard}}$ and the gradient is that of $y_{\text{soft}}$.
  \item \textbf{Gumbel-softmax}~\citep{c10-jang2017gumbel,c10-maddison2017concrete}. To sample a
        category from probabilities $\pi_1,\dots,\pi_K$, add independent Gumbel noise
        $\gamma_i=-\log(-\log u_i)$, $u_i\sim\mathrm{U}(0,1)$, and take $\argmax_i(\log\pi_i+\gamma_i)$.
        This is an exact sample. The relaxation replaces the $\argmax$ by
        $\softmax((\log\pi+\vec{\gamma})/\tau_g)$, which is differentiable and tends to a one-hot vector
        as $\tau_g\to 0$. The ``hard'' variant combines it with the straight-through trick. CommVQ
        trains its value encoder this way.
  \item \textbf{No gradient needed.} If the score is a fixed function of quantities that are trained
        anyway (an entropy or a prediction error) and $\tau$ is a hyperparameter, nothing about the
        gate itself needs to be learned. Only the threshold has to be chosen.
\end{enumerate}

% ===========================================================================
\section{AMOR: entropy-gated attention refinement}\label{sec:c13-amor}
% ===========================================================================

\subsection{Motivation}

A recurrent language model such as Mamba-2 or Gated DeltaNet (\cref{ch:linear,ch:delta}) is cheap
and handles most tokens well, but its fixed-size state cannot recall an arbitrary detail exactly
(\cref{ch:hybrid}). Fixed hybrids fix this by putting attention in some layers for \emph{every}
token. AMOR~\citep{zheng2026amor}, the Adaptive Metacognitive Output Router, takes the view of
dual-process theories of cognition~\citep{c13-kahneman2011}: a fast system handles routine input, and
a slow, deliberate system is engaged only when the fast one signals difficulty. The fast system is a
complete recurrent language model. The slow system is a small stack of causal attention blocks placed
\emph{after} the backbone. The ``metacognitive'' signal is the normalized entropy of the model's own
next-token prediction (\cref{sec:c13-entropy}). In the survey's words, AMOR appends ``post-hoc attention
refinement blocks to a recurrent backbone'' and activates them ``only for positions with high
normalized output entropy''~\citep{zhoubian2026memory}.

\begin{pitfall}[title={Which version of AMOR? Provenance of the code}]
The mirrored repository \repo{HaoranZhengRaul__AMOR} is on the first author's GitHub account, but
its README does not cite the paper and may follow a revised version. An earlier arXiv version of the
paper was titled ``Entropy-Based Metacognitive Gate for Dynamic SSM-Attention Switching'' and
described ``Ghost KV'': keys and values projected from the SSM \emph{hidden states}, reusing the
backbone's linear-time computation. The README states instead that ``K, V are standard linear
projections of the residual stream --- not SSM internal state'' (\repofile{repos/HaoranZhengRaul__AMOR/README.md}).
This chapter follows the code. The gating mechanism is the same in both descriptions.
\end{pitfall}

\begin{taxonomy}
\textbf{Representation:} implicit; the memory is the backbone's recurrent state plus the KV
caches of the AMOR attention blocks. \textbf{Update dynamics:} online. \textbf{Persistence:} long-term
(survey Table~I lists AMOR as ``Entropy-Gated Attention Refinement'', online, long-term).
\textbf{Update rule:} signal-gated routing (survey Table~II): an uncertainty signal decides whether the
expensive memory path is activated~\citep{zhoubian2026memory}. Note that AMOR gates the \emph{read}
(whether attention runs), not the write: as we will see, every token's key and value is still stored.
\end{taxonomy}

\subsection{Formulation}\label{sec:c13-amor-form}

\begin{figure}[t]
\centering
\begin{tikzpicture}[font=\small,>=Stealth,
  box/.style={draw,rounded corners=2pt,minimum height=7mm,align=center,inner sep=3pt},
  op/.style={draw,circle,inner sep=1pt,minimum size=5mm}]
\node[box] (bb) at (0,0.3) {recurrent\\backbone};
\node[box] (hin) at (2.0,0.3) {$\hat{\vh}_t$};
\node[box] (lm) at (4.2,1.3) {LM head $\mW_E$\\(no gradient)};
\node[box] (ent) at (6.8,1.3) {$\bar{\mathrm{H}}_t$};
\node[box] (g) at (9.3,1.3) {$g_t=\mathbf{1}[\bar{\mathrm{H}}_t>\tau]$};
\node[box] (att) at (5.0,-0.6) {causal attention\\$\mQ,\mK,\mV$ from $\hat{\vh}$, RoPE};
\node[op] (mul) at (9.3,-0.6) {$\times$};
\node[op] (add) at (10.6,-0.6) {$+$};
\node[box] (out) at (12.4,-0.6) {next block\\or LM head};
\draw[->] (bb) -- (hin);
\draw[->] (hin) |- (lm);
\draw[->] (hin) |- (att);
\draw[->] (lm) -- (ent);
\draw[->] (ent) -- (g);
\draw[->] (g) -- (mul);
\draw[->] (att) -- (mul);
\draw[->] (mul) -- node[above]{$\valpha\had$} (add);
\draw[->] (add) -- (out);
\draw[->] (hin.south) -- (2.0,-1.7) -| (add.south);
\end{tikzpicture}
\caption{One AMOR block. The block reads the normalized residual stream, measures the normalized
entropy of the prediction the LM head would make from it, and adds an attention update only where
the entropy exceeds a running threshold. $K$ such blocks are stacked after the backbone.}
\label{fig:c13-amor}
\end{figure}

\paragraph{Backbone.} Token embeddings pass through $N$ backbone blocks, each a recurrent mixer
(Mamba-2 or Gated DeltaNet) followed by a pre-norm SwiGLU MLP. Call the result $\vh_t\in\R^d$. In the
default (\code{classic}) residual mode it is normalized once, $\vh^{(0)}_t=\RMSNorm_f(\vh_t)$.

\paragraph{The entropy gate.} AMOR block $i\in\{1,\dots,K\}$ reads a normalized stream
$\hat{\vh}^{(i)}_t$ (equal to $\vh^{(0)}_t$ for $i=1$ and to $\RMSNorm_i(\vh^{(i-1)}_t)$ for $i>1$). It
applies the tied LM head $\mW_E\in\R^{\abs{\mathcal{V}}\times d}$ (the embedding matrix) and measures
the normalized entropy of the resulting distribution:
\begin{equation}\label{eq:c13-amor-entropy}
  \vz^{(i)}_t = \mW_E\,\hat{\vh}^{(i)}_t,\qquad
  \bar{\mathrm{H}}^{(i)}_t = \frac{\mathrm{H}\bigl(\softmax(\vz^{(i)}_t)\bigr)}{\log\abs{\mathcal{V}}}\in[0,1].
\end{equation}
This is the ``logit lens'': the prediction the model would make if it stopped at this point. The
gate compares it with a threshold that tracks the entropy distribution during training,
\begin{align}
  \bar\mu^{(i)} &\leftarrow m\,\bar\mu^{(i)} + (1-m)\,\operatorname{median}_{\text{batch}}\bar{\mathrm{H}}^{(i)},\qquad
  \bar\sigma^{(i)} \leftarrow m\,\bar\sigma^{(i)} + (1-m)\,\operatorname{std}_{\text{batch}}\bar{\mathrm{H}}^{(i)},\label{eq:c13-amor-ema}\\
  \tau^{(i)} &= \bar\mu^{(i)} + \kappa\,\bar\sigma^{(i)},\qquad
  g^{(i)}_t = \mathbf{1}\bigl[\bar{\mathrm{H}}^{(i)}_t>\tau^{(i)}\bigr],\label{eq:c13-amor-gate}
\end{align}
with momentum $m=0.99$, $\kappa=0.2$ and initial values $\bar\mu=0.3$, $\bar\sigma=0.1$ (the
\code{adaptive} mode; a \code{fixed} mode uses $\tau=\bar\mu+0.03$). The EMAs are updated only in
training and frozen at inference. The gate has no learnable parameters and is \emph{detached}: no
gradient flows through \eqref{eq:c13-amor-entropy} (\cref{sec:c13-hard}, option~1). By
\eqref{eq:c13-fire-rate}, the adaptive offset aims at a fire rate that does not depend on the scale of the
entropies, which change across model sizes and during training.

\paragraph{The refinement.} Where the gate fires, a causal multi-head attention over the block's own
keys and values refines the stream. For head $n$ with head dimension $d_h$,
\begin{equation}\label{eq:c13-amor-attn}
  \vq_t = \bm{R}_t\mW_Q^{n}\hat{\vh}_t,\quad \vk_s = \bm{R}_s\mW_K^{n}\hat{\vh}_s,\quad \vv_s=\mW_V^{n}\hat{\vh}_s,\qquad
  \vec{a}^{n}_t = \sum_{s\le t}\frac{\exp(\vq_t\T\vk_s/\sqrt{d_h})}{\sum_{s'\le t}\exp(\vq_t\T\vk_{s'}/\sqrt{d_h})}\,\vv_s,
\end{equation}
where $\bm{R}_t$ is the RoPE rotation (\cref{sec:c13-rope}) and the block index $i$ is dropped. The heads
are concatenated and projected by $\mW_O$, and the stream is updated with a per-channel weight:
\begin{equation}\label{eq:c13-amor-update}
  \vh^{(i)}_t = \vh^{(i-1)}_t + \valpha^{(i)}\had\bigl(g^{(i)}_t\,\mW_O[\vec{a}^1_t;\dots;\vec{a}^{H}_t]\bigr),
  \qquad \valpha^{(i)}=\sigmoid\bigl(\max(\vr^{(i)},-1)\bigr).
\end{equation}
$\mW_O$ starts at zero, so at initialization every AMOR block is the identity and the model is
exactly the backbone language model. The raw weights $\vr$ start at $0$ ($\valpha=0.5$), and the clamp
keeps each channel's $\alpha$ above $\sigmoid(-1)\approx 0.27$. The final logits are
$\mW_E\RMSNorm_{\text{out}}(\vh^{(K)}_t)$.

\paragraph{Stacking.} Each block recomputes its own entropy from the refined stream. Positions
resolved by block~1 drop in entropy and stop firing in block~2, while positions that stay uncertain
keep firing. Each block has its own EMAs, projections, RoPE and $\valpha$.

\paragraph{Training.} The loss is the usual next-token cross-entropy on the final logits. Because
$g_t$ is a constant for backpropagation, the attention parameters receive gradient only through the
outputs of firing positions. Non-firing positions still matter: their keys and values are attended
to by firing queries, and receive gradient that way. The backbone is trained ``exclusively through the
final logits path, never through the gate'' (README).

\paragraph{Dense-masked and sparse forms are equal.} Training multiplies a \emph{dense} causal
attention by the gate, which is GPU-friendly. Inference computes attention only at firing positions.
These agree.

\begin{proposition}\label{prop:c13-amor-sparse}
Let $\mA=\softmax(\mQ\mK\T/\sqrt{d_h}+\mM)\mV$ be causal attention and $\vg\in\{0,1\}^T$. Then
$\diag(\vg)\mA$ equals the matrix whose row $t$ is $\softmax(\mK_{1:t}\vq_t/\sqrt{d_h})\T\mV_{1:t}$ if
$g_t=1$ and $\bm{0}$ otherwise.
\end{proposition}
\begin{proof}
Row $t$ of the score matrix is $\vq_t\T\mK\T/\sqrt{d_h}$ plus row $t$ of $\mM$, which is $-\infty$ for
$s>t$. The softmax is taken row by row, so row $t$ of $\softmax(\cdot)$ is
$\softmax(\mK_{1:t}\vq_t/\sqrt{d_h})$ padded with zeros, and row $t$ of $\mA$ depends only on
$\vq_t,\mK_{1:t},\mV_{1:t}$. Multiplying on the left by $\diag(\vg)$ scales row $t$ by $g_t$. So rows
with $g_t=0$ become zero and rows with $g_t=1$ are unchanged.
\end{proof}

The proof also shows what the sparse form must keep: row $t$ needs \emph{all} keys and values up to
$t$, including those of positions that did not fire. This is why the released decoder stores every
token's key and value (\cref{alg:c13-amor-decode}).

\subsection{Algorithms and complexity}\label{sec:c13-amor-alg}

\begin{algorithm}[t]
\caption{AMOR forward pass (training and prefill), \texttt{classic} residual mode}
\label{alg:c13-amor-forward}
\begin{algorithmic}[1]
\Require tokens $x_{1:T}$; backbone; $K$ AMOR blocks; tied LM head $\mW_E$
\State $\mH\gets\RMSNorm_f(\text{Backbone}(x_{1:T}))\in\R^{T\times d}$
\For{$i=1,\dots,K$}
  \State $\hat{\mH}\gets\mH$ \textbf{if} $i=1$ \textbf{else} $\RMSNorm_i(\mH)$
  \State $\bar{\mathrm{H}}_t\gets$ normalized entropy of $\softmax(\mW_E\hat{\vh}_t)$ for all $t$ \Comment{no gradient}
  \If{training}
    \State update $\bar\mu^{(i)},\bar\sigma^{(i)}$ by \eqref{eq:c13-amor-ema}
  \EndIf
  \State $g_t\gets\mathbf{1}[\bar{\mathrm{H}}_t>\bar\mu^{(i)}+\kappa\bar\sigma^{(i)}]$ for all $t$
  \State $\mA\gets\mW_O\,\mathrm{CausalAttn}(\mathrm{RoPE}(\hat{\mH}\mW_Q),\mathrm{RoPE}(\hat{\mH}\mW_K),\hat{\mH}\mW_V)$ \Comment{dense}
  \State $\mH\gets\mH+\diag(\vg)\,\mA\,\diag(\valpha^{(i)})$
\EndFor
\State \Return logits $\RMSNorm_{\text{out}}(\mH)\,\mW_E\T$
\end{algorithmic}
\end{algorithm}

\begin{algorithm}[t]
\caption{AMOR decoding step at position $t$}
\label{alg:c13-amor-decode}
\begin{algorithmic}[1]
\Require token $x_t$; backbone recurrent states; per-block caches $(\mK^{(i)},\mV^{(i)})$
\State $\vh\gets\RMSNorm_f(\text{BackboneStep}(x_t))$ \Comment{updates the $\bigO(1)$ recurrent states}
\For{$i=1,\dots,K$}
  \State $\hat{\vh}\gets\vh$ \textbf{if} $i=1$ \textbf{else} $\RMSNorm_i(\vh)$
  \State $e\gets$ entropy of $\softmax(\mW_E\hat{\vh})$, or the fitted router $r_i(\hat{\vh})$
  \State append $\bm{R}_t\mW_K\hat{\vh}$ to $\mK^{(i)}$ and $\mW_V\hat{\vh}$ to $\mV^{(i)}$ \Comment{always}
  \If{$e>\tau^{(i)}$}
    \State $\vh\gets\vh+\valpha^{(i)}\had\mW_O\,\mathrm{Attn}(\bm{R}_t\mW_Q\hat{\vh},\mK^{(i)},\mV^{(i)})$
  \EndIf
\EndFor
\State \Return logits $\mW_E\RMSNorm_{\text{out}}(\vh)$
\end{algorithmic}
\end{algorithm}

\paragraph{Cost.} Let $f$ be the fraction of positions that fire. Per AMOR block, dense-masked training
costs $\bigO(T^2d)$ for attention plus $\bigO(Td\abs{\mathcal{V}})$ for the gate's LM-head product (with
$K=3$, the output layer is computed four times per position). Sparse inference costs $\bigO(Td^2)$ for the
key and value projections of all positions and $\bigO(fT^2d)$ for the queries and attention of firing
positions. The decoding memory is $2td$ numbers per block after $t$ tokens, independent of $f$: for the 180M
configuration ($d=768$, $K=3$), $3\cdot2\cdot768=4608$ numbers per token, against $12\cdot2\cdot768=18432$ for
a twelve-layer Transformer of the same width. So AMOR reduces attention \emph{compute} in proportion to
the fire rate, and attention \emph{memory} only through the architectural choice of few attention blocks
($K\ll L$), not through the gate.

\subsection{Implementation walkthrough}\label{sec:c13-amor-impl}

\begin{implnote}[title={In the code: map of \texttt{repos/HaoranZhengRaul\_\_AMOR}}]
\begin{itemize}
  \item \repofile{repos/HaoranZhengRaul__AMOR/amor.py}: building blocks (\code{RMSNorm},
        \code{SwiGLUMLP}, \code{RotaryPositionEmbedding}, \code{KVProjection}), the gate
        \code{EntropyGate}, the block \code{AMORBlock}, the Gated DeltaNet backbone block
        \code{GDNBlock}, and the unified model \code{AMOR}.
  \item \repofile{repos/HaoranZhengRaul__AMOR/amor_decode.py}: \code{AMORDecodeWrapper} with
        \code{prefill} and \code{decode_step}, and \code{_chunked_gate}.
  \item \repofile{repos/HaoranZhengRaul__AMOR/amor_router.py}: the fitted \code{EntropyRouter},
        \code{attach_routers} and \code{detach_routers}. \code{load_model.py}, \code{generate.py}
        and \code{weight_io.py} handle checksum-verified loading and greedy generation.
\end{itemize}
\textbf{Variants.} Six variants are flag combinations of one class: \code{backbone} in
\{\code{'mamba2'}, \code{'gdn'}\}, \code{n_amor_blocks} in \{1, 3\} and \code{residual_mode} in
\{\code{'classic'}, \code{'h'}\}. In \code{'h'} mode the backbone's final norm is used as block~1's
pre-norm and the residual stream stays un-normalized until \code{final_norm}.
\textbf{Backbones.} For Mamba-2 the code builds a \code{transformers} \code{Mamba2ForCausalLM},
keeps its embeddings, SSD layers and \code{norm_f}, and adds a separate \code{RMSNorm} and
\code{SwiGLUMLP} after each SSD layer. For GDN it wraps \code{fla.layers.GatedDeltaNet} with
\texttt{num\_heads = int(0.75 * d\_model / head\_dim)} and \code{expand_v=2.0}.
\textbf{Defaults} (\code{AMOR.__init__}): \code{vocab_size=128256}, \code{d_model=768},
\code{n_layer=12}, \code{d_ff=1216}, \code{head_dim=64}, \code{n_heads=12}, \code{d_state=128},
\code{init_threshold=0.3}, \code{gate_offset=0.03}, \code{offset_mode='adaptive'},
\code{offset_k=0.2}, \code{attention_mode='full_with_mask'}, \code{max_seq_len=4096}. All weights
are initialized from $\mathcal{N}(0,0.02^2)$, \code{out_proj} stays zero, and the LM head is tied to
the embedding. The README lists 180M ($d=768$, 12 layers), 440M ($d=1024$, 24 layers) and 1.5B
($d=2048$, 24 layers) scales, with a Llama-3.1-compatible tokenizer.
\end{implnote}

The gate is short. The excerpt omits the branch that calls a fitted router.

\begin{lstlisting}[caption={Excerpt from \texttt{repos/HaoranZhengRaul\_\_AMOR/amor.py} (\texttt{EntropyGate.forward})},label={lst:c13-amor-gate}]
    def forward(self, logits: torch.Tensor) -> Tuple[torch.Tensor, torch.Tensor]:
        # ...
        else:
            logits_f32 = logits.detach().float()
            probs = F.softmax(logits_f32, dim=-1)
            raw_entropy = -(probs * (probs + 1e-10).log()).sum(dim=-1)
            entropy = raw_entropy / self.max_entropy

        if self.training:
            entropy_detached = entropy.detach()
            batch_median = entropy_detached.median()
            self.running_threshold.mul_(self.ema_momentum).add_(
                batch_median, alpha=1.0 - self.ema_momentum
            )
            if self.offset_mode == 'adaptive':
                batch_std = entropy_detached.std()
                self.running_std.mul_(self.ema_momentum).add_(
                    batch_std, alpha=1.0 - self.ema_momentum
                )

        if self.offset_mode == 'adaptive':
            threshold = self.running_threshold + self.offset_k * self.running_std
        else:
            threshold = self.running_threshold + self.offset

        gate = (entropy > threshold).float()
        return gate, entropy
\end{lstlisting}

\begin{implnote}[title={In the code: dense-masked and true-sparse attention}]
\code{AMORBlock} has two attention paths. \code{_attention_full_with_mask} projects $\mQ,\mK,\mV$ to
shape \texttt{[batch, n\_heads, seq, head\_dim]}, applies RoPE, calls
\texttt{F.scaled\_dot\_product\_attention(Q, K, V, is\_causal=True)}, applies \code{out_proj} and multiplies
by \code{gate.unsqueeze(-1)}: the left side of \cref{prop:c13-amor-sparse}.
\code{_attention_true_sparse} computes $\mK,\mV$ for all positions, then gathers the firing rows with
\texttt{active\_mask = gate > 0.5}, computes $\mQ$ only for them (with RoPE at their own positions,
\code{cos_cached[seq_indices]}), builds the causal mask by comparing \texttt{key\_positions > positions},
and scatters the outputs back: the right side. Its docstring calls the two ``bit-identical'', while the
README more carefully says that ``small floating-point differences from dense attention can occur''
(different kernels sum in different orders). The entropy itself is computed in fp32 with
$\epsilon=10^{-10}$ inside the log, which is the direct form \eqref{eq:c13-entropy} rather than the
log-sum-exp form \eqref{eq:c13-entropy-lse}.
\end{implnote}

The decoder follows \cref{alg:c13-amor-decode}. Note the comment in the excerpt's source file header:
``K and V are projected and appended unconditionally so that the cache stays causal regardless of
past gate decisions.''

\begin{lstlisting}[caption={Excerpt from \texttt{repos/HaoranZhengRaul\_\_AMOR/amor\_decode.py} (\texttt{decode\_step})},label={lst:c13-amor-decode}]
            fired = gate.item() > 0.5

            K_new, V_new = ab.kv_proj(h_normed)
            K_new = K_new.view(batch, 1, n_heads, head_dim).transpose(1, 2)
            V_new = V_new.view(batch, 1, n_heads, head_dim).transpose(1, 2)
            cos = ab.rope.cos_cached[step_idx:step_idx + 1].unsqueeze(0).unsqueeze(0)
            sin = ab.rope.sin_cached[step_idx:step_idx + 1].unsqueeze(0).unsqueeze(0)
            K_new = RotaryPositionEmbedding._apply_rotary(K_new, cos, sin)

            K_full = torch.cat([past_K, K_new], dim=2)
            V_full = torch.cat([past_V, V_new], dim=2)
            new_amor_kv_caches.append((K_full, V_full))

            if fired:
                Q = ab.q_proj(h_normed).view(batch, 1, n_heads, head_dim).transpose(1, 2)
                Q = RotaryPositionEmbedding._apply_rotary(Q, cos, sin)
                attn_out = F.scaled_dot_product_attention(Q, K_full, V_full, is_causal=False)
                attn_out = attn_out.transpose(1, 2).contiguous().view(batch, 1, ab.d_model)
                attn_out = ab.out_proj(attn_out)
                alpha = ab.get_alpha()
                h = self._residual_step(h, h_normed, alpha * attn_out, pre_norm)
\end{lstlisting}

\paragraph{The fitted router.}\label{sec:c13-amor-router} Each gate costs an LM-head product
(\cref{sec:c13-entropy}). For inference, the repository can replace it by a small network that predicts
the normalized entropy directly from the same normalized residual,
\begin{equation}\label{eq:c13-amor-router}
  r(\hat{\vh}) = \vw_2\T\SiLU(\mW_1\hat{\vh}+\vb_1)+b_2+\vw_s\T\hat{\vh},\qquad \mW_1\in\R^{512\times d},
\end{equation}
a width-512 SiLU MLP with a linear skip, compared with the \emph{same} frozen threshold. Its cost is
about $512d+512+d$ multiply-adds, roughly $0.39\times 10^6$ for $d=768$, against $9.85\times10^{7}$ for
the LM head: a factor of about $250$.

\begin{implnote}[title={In the code: routers and chunked gates}]
\code{EntropyRouter} (\repofile{repos/HaoranZhengRaul__AMOR/amor_router.py}) keeps \code{fc1},
\code{fc2} and \code{skip} in fp32 even when the model is cast to bf16 (it overrides \code{_apply} and
disables autocast). \code{attach_routers} accepts only an eval-mode model whose base-weight SHA-256 matches
the router package (\code{variant} \code{"distill"}), and \code{EntropyGate.forward} refuses a router in
training mode. Without routers, \code{_chunked_gate} evaluates the LM head over the prompt in chunks of
4096 positions so that the full \texttt{[B, L, V]} logit tensor is never materialized.
\end{implnote}

\subsection{Reference implementation}

\Cref{lst:c13-amor-ref} is a self-contained version of one AMOR block. It keeps the essential
choices: the detached normalized entropy, the EMA median-plus-std threshold, dense causal attention
multiplied by the gate, a zero-initialized output projection and the clamped per-channel $\valpha$.

\begin{lstlisting}[caption={Reference implementation (ours): entropy-gated attention refinement},label={lst:c13-amor-ref}]
import math, torch, torch.nn as nn, torch.nn.functional as F

class EntropyGatedRefinement(nn.Module):
    def __init__(self, d, n_heads, lm_head, momentum=0.99, k_off=0.2):
        super().__init__()
        self.h, self.dh = n_heads, d // n_heads
        self.qkv = nn.Linear(d, 3 * d, bias=False)
        self.out = nn.Linear(d, d, bias=False)
        nn.init.zeros_(self.out.weight)            # block starts as the identity
        self.raw_alpha = nn.Parameter(torch.zeros(d))
        self.lm_head = lm_head                     # tied output layer, shared
        self.register_buffer("mu", torch.tensor(0.3))
        self.register_buffer("sd", torch.tensor(0.1))
        self.m, self.k_off = momentum, k_off

    @torch.no_grad()
    def gate(self, h):                             # h: [B, T, d], normalized
        logp = torch.log_softmax(self.lm_head(h).float(), dim=-1)
        ent = -(logp.exp() * logp).sum(-1) / math.log(logp.shape[-1])
        if self.training:                          # track median and std
            self.mu.mul_(self.m).add_((1 - self.m) * ent.median())
            self.sd.mul_(self.m).add_((1 - self.m) * ent.std())
        return (ent > self.mu + self.k_off * self.sd).float(), ent

    def forward(self, h, rope):                    # rope: [B,H,T,dh] -> same
        B, T, d = h.shape
        g, ent = self.gate(h)
        q, k, v = self.qkv(h).view(B, T, 3, self.h, self.dh).permute(2, 0, 3, 1, 4)
        a = F.scaled_dot_product_attention(rope(q), rope(k), v, is_causal=True)
        a = self.out(a.transpose(1, 2).reshape(B, T, d))
        alpha = torch.sigmoid(self.raw_alpha.clamp(min=-1.0))
        return h + alpha * a * g.unsqueeze(-1), g, ent
\end{lstlisting}

We checked \cref{prop:c13-amor-sparse} numerically in NumPy (dense masked attention against attention
computed row by row at firing positions; maximum difference $10^{-16}$), and the bounds
$0\le\bar{\mathrm{H}}\le 1$ on random logits.

\begin{results}
\textbf{Source:} \repofile{repos/HaoranZhengRaul__AMOR/README.md}. The repository contains no
numerical results. It states that trained weights ``will be linked here after the review period'', and
describes the design goal qualitatively: ``The entropy gate concentrates attention compute on positions
where the backbone is least confident.'' The survey describes AMOR in the same terms and does not quote
numbers~\citep{zhoubian2026memory}. The paper reports experiments on synthetic retrieval tasks and
language modelling, comparing against recurrent-only and Transformer-only baselines; see the paper for
numbers.
\end{results}

\subsection{Discussion}

\paragraph{Strengths.} AMOR leaves the backbone a complete language model and adds refinement as an
option, so it degrades gracefully: at initialization, and wherever the gate is closed, the model is
the backbone. The signal is interpretable (one can inspect where attention fired), and the compute it
spends follows the model's own uncertainty.

\paragraph{Limitations.} (i) In the released decoder, the KV memory of the attention blocks still grows
with every token, so the gate saves compute, not memory. (ii) The native gate needs an LM-head product per
block, which is why the fitted router exists. (iii) Entropy conflates reducible and irreducible
uncertainty (\cref{sec:c13-entropy}). (iv) The threshold is calibrated on training batches; a shift in
the input distribution at inference shifts the fire rate. (v) Attention sits only at the top, so the
backbone's own layers never see retrieved information. Compare the fixed hybrids of \cref{ch:hybrid},
which interleave attention throughout the depth.

\paragraph{Relations.} AMOR is a form of conditional computation, related to adaptive computation
time~\citep{c13-graves2016act} and to Mixture-of-Depths, which routes a fixed fraction of tokens through
each block~\citep{c13-raposo2024mod}. Unlike these, its router is not learned but derived from the
model's output distribution. HAM, next, moves the decision from the read to the write.

% ===========================================================================
\section{HAM: Hybrid Associative Memories}\label{sec:c13-ham}
% ===========================================================================

\subsection{Motivation}

A recurrent layer and an attention layer store the past in opposite ways. The recurrent state
compresses \emph{everything} into $d_k\times d_v$ numbers. The KV cache stores \emph{everything}
verbatim. HAM~\citep{lufkin2026ham} (Zyphra) combines them inside one sequence-mixing layer by a
division of labour. The authors state it as follows: the RNN compresses the entire sequence, while
attention supplements it ``only with information that is difficult for the RNN to predict, which is
hence the most valuable information to explicitly store.'' The survey describes the layer as ``a
recurrent memory [that] processes every token, whereas a sparse KV cache stores only tokens that are
difficult for the recurrent state to predict, with the cache growth controlled by a threshold or
learned router''~\citep{zhoubian2026memory}. The cache therefore grows with how much \emph{novel}
information the sequence contains, not with its length, and a single continuous threshold controls the
trade-off between cache size and quality.

The survey places HAM next to the hybrid quadratic-linear Transformer of
\citet{irie2025blending} (\cref{ch:hybrid}), a ``fixed-budget precursor'' that blends a fast-weight
memory with a bounded KV memory. HAM differs ``by making cache admission data-dependent and
adjustable at runtime.'' The idea has a long history in neuroscience: complementary learning systems
theory assigns fast, verbatim storage of surprising episodes to the hippocampus and slow statistical
learning to the neocortex~\citep{c13-mcclelland1995cls}.

\begin{taxonomy}
\textbf{Representation:} implicit (recurrent state plus KV cache). \textbf{Update dynamics:}
online. \textbf{Persistence:} long-term (survey Table~I: ``Surprise-Gated KV-RNN Hybrid Memory'').
\textbf{Update rule:} HAM appears in two rows of survey Table~II, signal-gated writing and
admission/eviction, because a prediction-error signal decides admission into the cache~\citep{zhoubian2026memory}.
\end{taxonomy}

\subsection{Formulation}\label{sec:c13-ham-form}

There is no public code, and the survey and abstract do not give equations. \emph{One way to
formalize the layer is the following; the paper's exact parameterization may differ.}

\paragraph{Recurrent memory.} Any recurrent memory from \cref{ch:linear,ch:delta} can serve. We use
the gated delta rule, whose write is itself driven by prediction error:
\begin{equation}\label{eq:c13-ham-rnn}
  \mS_t = \alpha_t\mS_{t-1}\bigl(\mI-\beta_t\vk_t\vk_t\T\bigr)+\beta_t\vv_t\vk_t\T,
  \qquad \norm{\vk_t}=1 .
\end{equation}

\paragraph{Surprise.} Before writing token $t$, measure how well the memory already recalls its value,
using the normalized prediction error \eqref{eq:c13-normerr}:
\begin{equation}\label{eq:c13-ham-surprise}
  s_t = \frac{\norm{\vv_t-\mS_{t-1}\vk_t}^2}{\norm{\vv_t}^2+\epsilon}.
\end{equation}
A learned alternative is a router $s_t=\sigmoid(\vw\T\vx_t+b)$ on the layer input, trained with a
straight-through estimator or an auxiliary sparsity loss (\cref{sec:c13-hard}). The survey mentions
both a threshold and a learned router; we develop the threshold form.

\paragraph{Admission and cache.} With threshold $\tau$,
\begin{equation}\label{eq:c13-ham-admit}
  g_t=\mathbf{1}[s_t>\tau],\qquad
  \mathcal{C}_t=\begin{cases}\mathcal{C}_{t-1}\cup\{t\} & g_t=1,\\ \mathcal{C}_{t-1} & g_t=0.\end{cases}
\end{equation}

\paragraph{Read-out.} The output adds an attention read over the admitted entries to the recurrent
read-out:
\begin{equation}\label{eq:c13-ham-read}
  \vo_t = \mS_t\vq_t + \lambda_t\sum_{s\in\mathcal{C}_t}
  \frac{\exp(\vq_t\T\vk_s/\sqrt{d_k})}{\sum_{s'\in\mathcal{C}_t}\exp(\vq_t\T\vk_{s'}/\sqrt{d_k})}\,\vv_s ,
\end{equation}
with the attention term set to $\bm{0}$ while $\mathcal{C}_t$ is empty, and $\lambda_t$ a learned
(possibly data-dependent) mixing weight. Other choices are possible, for example storing the
innovation $\vv_s-\mS_{s-1}\vk_s$ instead of $\vv_s$, so that attention supplies exactly what the
RNN misses.

\paragraph{Why prediction error is the right signal.} Combine \eqref{eq:c13-ham-rnn} and
\eqref{eq:c13-ham-surprise}. Take $\alpha_t=\beta_t=1$ and orthonormal keys. After writing
$(\vk_t,\vv_t)$, the memory recalls $\mS_t\vk_t = \mS_{t-1}\vk_t-\mS_{t-1}\vk_t\vk_t\T\vk_t+\vv_t\vk_t\T\vk_t=\vv_t$
exactly, and writes along other orthonormal keys leave this recall unchanged. So a pair that occurs again has
$s=0$ and is never admitted twice. The cache \emph{deduplicates} what the RNN has already learned. When
the RNN runs out of capacity (more than $d_k$ distinct keys, decay $\alpha_t<1$, or interference between
non-orthogonal keys), recall degrades, $s_t$ rises, and more tokens are admitted. The cache grows exactly
where the fixed-size state fails.

\begin{example}[Deduplication]\label{ex:c13-ham-dedup}
Take $d_k=d_v=8$, six orthonormal keys with random values, and the stream of item indices
$0,1,2,0,3,1,4,0,5,2$, with $\alpha=\beta=1$ and $\tau=0.5$. The first occurrence of each item has
$s_t=1$ (the memory recalls $\bm{0}$), and every repeat has $s_t=0$. The admitted positions are
$\{0,1,2,4,6,8\}$, one per distinct item. We checked this in NumPy, together with the fact that the
cache size is non-increasing in $\tau$ on a stream with decay and partial repetition.
\end{example}

\paragraph{Cache growth.} By \eqref{eq:c13-expected-cache}, $\E\abs{\mathcal{C}_T}=\sum_t\Pr(s_t>\tau)$.
The two extremes are $\tau<0$, which admits every token and gives an RNN plus full attention, and
$\tau\ge\max_t s_t$, which admits nothing and gives a pure RNN. Because $\tau$ enters only the
admission test, it can be changed at inference time without retraining. How robust a trained model is
to such changes depends on training (one option is to sample $\tau$ per training sequence); see the
paper for how this is handled.

\paragraph{Parallel training form.} Recurrence \eqref{eq:c13-ham-rnn} is trained with the chunkwise
algorithms of \cref{ch:delta}. For the delta rule these already compute the pseudo-values
$\vu_t=\beta_t(\vv_t-\mS_{t-1}\vk_t)$ for all $t$ in parallel (the WY/UT form), so the surprise
$s_t=\norm{\vu_t}^2/(\beta_t^2\norm{\vv_t}^2)$ costs nothing extra in the ungated case. Given the mask
$\vg$, the attention term of \eqref{eq:c13-ham-read} for all $t$ is a standard masked attention:
\begin{equation}\label{eq:c13-ham-parallel}
  \mA = \softmax\bigl(\mQ\mK\T/\sqrt{d_k}+\mM+\mM_{\vg}\bigr)\mV,\qquad
  (\mM_{\vg})_{ts}=\begin{cases}0 & g_s=1,\\ -\infty & g_s=0,\end{cases}
\end{equation}
where rows with no admitted key are set to zero. The argument is the column version of
\cref{prop:c13-amor-sparse}: a key with score $-\infty$ receives weight $e^{-\infty}=0$ and drops out of
both numerator and denominator, so row $t$ is exactly the softmax over $\mathcal{C}_t$. During training
this costs as much as full attention. The saving is at inference, where only admitted entries are stored.

\subsection{Algorithm and complexity}

\begin{algorithm}[t]
\caption{HAM-style layer, recurrent form (our formalization)}\label{alg:c13-ham}
\begin{algorithmic}[1]
\Require $\vq_t,\vk_t\in\R^{d_k}$ ($\norm{\vk_t}=1$), $\vv_t\in\R^{d_v}$, gates $\alpha_t,\beta_t$, threshold $\tau$
\State $\mS_0\gets\bm{0}$, $\mathcal{C}_0\gets\emptyset$
\For{$t=1,\dots,T$}
  \State $s_t\gets\norm{\vv_t-\mS_{t-1}\vk_t}^2/(\norm{\vv_t}^2+\epsilon)$ \Comment{surprise}
  \If{$s_t>\tau$} \Comment{admit}
    \State append $(\vk_t,\vv_t)$ to the cache $\mathcal{C}$
  \EndIf
  \State $\mS_t\gets\alpha_t\mS_{t-1}(\mI-\beta_t\vk_t\vk_t\T)+\beta_t\vv_t\vk_t\T$ \Comment{every token}
  \State $\vo_t\gets\mS_t\vq_t+\lambda_t\,\mathrm{Attn}(\vq_t,\mathcal{C})$
\EndFor
\end{algorithmic}
\end{algorithm}

Per token, the RNN costs $\bigO(d_kd_v)$, the surprise another $\bigO(d_kd_v)$ (one matrix-vector
product), and the attention read $\bigO(\abs{\mathcal{C}_t}(d_k+d_v))$. The memory is
$d_kd_v+\abs{\mathcal{C}_t}(d_k+d_v)$ numbers per head. With admission rate $r$, a HAM layer stores
about $r$ times the KV cache of an attention layer. Unlike AMOR, the saving is in memory as well as in
compute.

\subsection{Reference implementation}

\begin{lstlisting}[caption={Reference implementation (ours): surprise-gated KV admission over a gated delta-rule memory},label={lst:c13-ham-ref}]
import torch, torch.nn.functional as F

def ham_layer(q, k, v, alpha, beta, tau, lam=1.0, eps=1e-6):
    """q, k: [T, d_k]; v: [T, d_v]; alpha, beta: [T] gates in (0, 1).
    Recurrent form of a surprise-gated KV-RNN hybrid (our formalization)."""
    T, d_k = k.shape
    S = k.new_zeros(v.shape[1], d_k)          # fixed-size recurrent memory
    cache_k, cache_v, admitted, out = [], [], [], []
    for t in range(T):
        kt = F.normalize(k[t], dim=-1)
        err = v[t] - S @ kt                     # what the memory fails to recall
        s = err.pow(2).sum() / (v[t].pow(2).sum() + eps)
        if s > tau:                             # admission: only surprising tokens
            cache_k.append(kt); cache_v.append(v[t]); admitted.append(t)
        # gated delta rule: S <- alpha * S (I - beta k k^T) + beta v k^T
        S = alpha[t] * (S - beta[t] * torch.outer(S @ kt, kt)) \
            + beta[t] * torch.outer(v[t], kt)
        o = S @ q[t]
        if cache_k:                             # exact recall from the sparse cache
            K, V = torch.stack(cache_k), torch.stack(cache_v)
            w = torch.softmax(K @ q[t] / d_k ** 0.5, dim=0)
            o = o + lam * (w @ V)
        out.append(o)
    return torch.stack(out), admitted
\end{lstlisting}

\begin{results}
\textbf{Source:} no public code was found (\repofile{README.md}, section~3), so there is no repository
file with results. The authors state publicly that the hybrid ``offers strong, competitive performance
relative to RNNs and Transformers even at substantially lower KV-cache usage'', and that controlling
the cache growth rate with the threshold gives ``a smooth trade-off with loss and performance''. The
survey makes the qualitative point that HAM makes ``cache admission data-dependent and adjustable at
runtime''~\citep{zhoubian2026memory}. See the paper for numbers.
\end{results}

\subsection{Discussion}

\paragraph{Strengths.} Memory grows with novelty rather than with length, which is the right scaling
for long inputs that are mostly predictable (code with repeated identifiers, long documents with
recurring entities). One continuous knob trades memory for quality at deployment time. The signal is
cheap and is a by-product of delta-rule training.

\paragraph{Limitations.} (i) Admission is irreversible: a token that was predictable when written but
becomes important later (a detail whose relevance is revealed by a later question) is lost to exact
recall. The survey lists ``delayed relevance'' among the risks of signal-gated writing. (ii) Surprise is
not utility: noise is surprising but useless, and a stream with many random tokens fills the cache.
(iii) The surprise depends on the RNN's state, so one threshold may admit very different rates in
different layers and heads. (iv) Different sequences in a batch admit different numbers of tokens, so
caches are ragged. Paged KV memory~\citep{kwon2023paged} (\cref{ch:foundations}) handles this naturally.

\paragraph{Relations.} AMOR gates the \emph{read} by output entropy, HAM gates the \emph{write} by
memory prediction error. Titans (\cref{ch:ttt}) uses a surprise signal to scale \emph{parametric}
writes, and HAM uses one to decide \emph{verbatim} storage. Expansion Span~\citep{nunez2025expansion}
(\cref{ch:hybrid}) also reserves attention capacity for selected far-away tokens, chosen by retrieval
rather than by surprise at write time.

% ===========================================================================
\section{Background II: managing the working memory}\label{sec:c13-mgmt-bg}
% ===========================================================================

Part B treats the working memory as a data structure to be engineered. This section builds the tools:
how large the KV cache is, how quantization shrinks it, why rotary embeddings complicate this, and
what it means to checkpoint a recurrent state.

\subsection{How large is the KV cache?}\label{sec:c13-kvsize}

A decoder with $L$ attention layers, $H_{kv}$ key-value heads of dimension $d_k$, and $b$ bytes per
number stores, per token,
\begin{equation}\label{eq:c13-kv-bytes}
  \text{bytes per token} = 2\cdot L\cdot H_{kv}\cdot d_k\cdot b ,
\end{equation}
the factor $2$ counting keys and values (\cref{ch:foundations}). The CommVQ scripts target
Llama-3.1-8B, with $L=32$ layers (\texttt{N = 32} in \repofile{repos/UMass-Embodied-AGI__CommVQ/training/make_scale.py})
and $H_{kv}d_k=8\cdot128=1024$ (\texttt{KV\_DIM = 1024} and the reshape \texttt{view(N, 8, 2, 64)} in
\repofile{repos/UMass-Embodied-AGI__CommVQ/training/quantize_key_cache.py}). In 16-bit precision,
\eqref{eq:c13-kv-bytes} gives $2\cdot32\cdot1024\cdot2=131072$ bytes, i.e.\ 128~KiB per token, or 16~GiB
for a single sequence of $2^{17}=131072$ tokens. That is the same order as the model's weights, and it
is multiplied by the batch size.

There are several independent levers, and the systems in this chapter pull different ones:
\emph{fewer or smaller heads} (grouped-query attention~\citep{c01-ainslie2023gqa}, already in the
$H_{kv}=8$ above, or low-rank latent caches~\citep{deepseek2024v2}); \emph{fewer tokens} (eviction and
admission, \cref{sec:c13-policies}; HAM); \emph{fewer bits per number} (quantization; CommVQ);
\emph{states instead of tokens} (recurrent layers; Memory Caching); \emph{better placement}
(PagedAttention~\citep{kwon2023paged}, \cref{ch:foundations}, stores the cache in fixed-size blocks
addressed through a block table, which removes fragmentation and lets sequences share prefix blocks; it
does not shrink the cache, so it combines with all the others); and \emph{better content} (rewriting
entries; Bottlenecked Transformers).

\subsection{Scalar quantization}\label{sec:c13-sq}

A $b$-bit \emph{uniform scalar quantizer} maps a real number $x$ to one of $2^b$ levels. With step
(scale) $\Delta$ and offset $z$,
\begin{equation}\label{eq:c13-sq}
  Q(x) = \Delta\cdot\bigl(\operatorname{clip}(\lfloor x/\Delta\rceil+z,\,0,\,2^b-1)-z\bigr),
\end{equation}
where $\lfloor\cdot\rceil$ rounds to the nearest integer. Inside the clipping range the error is at most
$\Delta/2$; if it is spread uniformly over $[-\Delta/2,\Delta/2]$ its mean square is
$\int_{-\Delta/2}^{\Delta/2}e^2\,de/\Delta=\Delta^2/12$. To cover a range $[x_{\min},x_{\max}]$ one sets
$\Delta=(x_{\max}-x_{\min})/(2^b-1)$, so a single outlier inflates the error of every other number in
its group. KV-cache quantizers therefore use small groups with their own scale and offset, and choose
the grouping per channel for keys (whose outliers sit in fixed channels) and per token for
values~\citep{c13-liu2024kivi}. The scales cost bits too. A group of $G$ numbers with a 16-bit scale
and a 16-bit offset costs $b+32/G$ bits per number. At $b=1$ or $2$ there are only two or four levels
per number, and the error is large whatever the grouping.

\subsection{Vector quantization, \texorpdfstring{$k$}{k}-means and EM}\label{sec:c13-vq}

A \emph{vector quantizer} replaces a whole vector $\vx\in\R^m$ by one of $K$ codewords
$\vec{c}_1,\dots,\vec{c}_K$ (the \emph{codebook}). The encoder stores only the index
$i(\vx)=\argmin_i\norm{\vx-\vec{c}_i}$, and the decoder looks up $\vec{c}_{i(\vx)}$~\citep{c13-gray1984vq}.
The rate is $\log_2K/m$ bits per coordinate. Scalar quantization of each coordinate with $b$ bits is the
special case whose $K=2^{bm}$ codewords form a rectangular grid. VQ can place the same number of
codewords anywhere, in particular where the data is and with shapes that follow correlations between
coordinates, so at equal rate its distortion is never worse and usually much better.

\paragraph{Lloyd's algorithm ($k$-means).} To fit a codebook to data $\vx_1,\dots,\vx_N$, minimize the
distortion $J=\sum_n\norm{\vx_n-\vec{c}_{i_n}}^2$ over codewords and assignments by alternating two
steps~\citep{c13-lloyd1982}. (1) With the codewords fixed, the best assignment of each $\vx_n$ is its
nearest codeword. (2) With the assignments fixed, $J$ is a sum of separate quadratics, one per codeword,
and setting $\partial J/\partial\vec{c}_i=-2\sum_{n:i_n=i}(\vx_n-\vec{c}_i)=\bm{0}$ gives the centroid
$\vec{c}_i=\text{mean}\{\vx_n:i_n=i\}$. Neither step can increase $J$, and there are finitely many
assignments, so the algorithm stops at a local minimum.

\paragraph{EM and annealing.} Lloyd's hard assignments are the zero-temperature limit of
expectation-maximization~\citep{c13-dempster1977} for a mixture of isotropic Gaussians with common
variance. The E-step computes soft responsibilities
\begin{equation}\label{eq:c13-estep}
  w_{nc} = \frac{\exp(-d(\vx_n,\vec{c}_c)/\tau)}{\sum_{c'}\exp(-d(\vx_n,\vec{c}_{c'})/\tau)} ,
\end{equation}
with $d$ the squared distance and $\tau=2\sigma^2$. The M-step sets each codeword to the
$w$-weighted mean. As $\tau\to0$ the softmax becomes the nearest-neighbour indicator and EM becomes
$k$-means. \emph{Deterministic annealing}~\citep{c13-rose1998da} starts at a high $\tau$, where all
codewords are pulled towards the data mean and poor local minima are smoothed away, and lowers $\tau$
gradually. CommVQ's key codebook is trained this way, with one change (\cref{sec:c13-commvq}): its
codewords are not free but constrained to a structured family.

\paragraph{The M-step for a structured codebook.} Suppose every codeword is a fixed linear function of
a parameter vector, $\vec{c}_c=\mT_c\vtheta$ (stack them as $\vec{c}=\mT\vtheta$). Let
$N_c=\sum_nw_{nc}$ and $\vm_c=N_c^{-1}\sum_nw_{nc}\vx_n$ be the soft count and weighted mean of
codeword $c$. Because $\sum_nw_{nc}(\vx_n-\vm_c)=\bm{0}$, the cross term vanishes in
\[
  \sum_n w_{nc}\norm{\vx_n-\vec{c}_c}^2 = \sum_n w_{nc}\norm{\vx_n-\vm_c}^2 + N_c\norm{\vm_c-\vec{c}_c}^2 ,
\]
so up to a constant the M-step objective is a weighted least-squares problem,
\[
  \sum_cN_c\norm{\vm_c-\mT_c\vtheta}^2=(\vm-\mT\vtheta)\T\mathbf{N}(\vm-\mT\vtheta),
\]
where $\mathbf{N}$ is diagonal with $N_c$ repeated over the coordinates of codeword $c$. Setting the gradient
$-2\mT\T\mathbf{N}(\vm-\mT\vtheta)$ to zero gives
\begin{equation}\label{eq:c13-structured-mstep}
  \vtheta^\star = (\mT\T\mathbf{N}\mT)^{-1}\mT\T\mathbf{N}\,\vm .
\end{equation}
Free codewords are the case $\mT=\mI$, which gives back $\vec{c}_c=\vm_c$.

\subsection{Large codebooks: product, residual and additive quantization}\label{sec:c13-rvq}

At 1 bit per coordinate a 128-dimensional head vector needs $K=2^{128}$ codewords, which cannot be
stored or searched. Three standard ways to get exponentially many codewords from a few small codebooks
are the following.
\begin{itemize}
  \item \textbf{Product quantization}~\citep{c13-jegou2011pq} splits $\vx$ into $M$ sub-vectors and
        quantizes each with its own codebook of size $K$. This gives $K^M$ codewords for $MK$ stored ones.
  \item \textbf{Residual VQ}~\citep{c13-gray1984vq} quantizes in stages. Stage $r$ quantizes what the
        earlier stages left over,
        \[
          \vec{\delta}^{(0)}=\vx,\qquad i_r=\argmin_i\norm{\vec{\delta}^{(r-1)}-\vec{c}^{(r)}_i},\qquad
          \vec{\delta}^{(r)}=\vec{\delta}^{(r-1)}-\vec{c}^{(r)}_{i_r},
        \]
        and decodes $\hat{\vx}=\sum_{r=1}^R\vec{c}^{(r)}_{i_r}$ with $R\log_2K$ bits. Each stage refines
        the previous ones. In a small NumPy test with unit vectors in $\R^4$ and three stages of 16
        codewords, the mean squared error fell by roughly a factor of three per stage.
  \item \textbf{Additive quantization}~\citep{c13-babenko2014aq} also decodes a sum of codewords, one per
        codebook, but chooses all indices jointly rather than greedily. Exact joint encoding is a hard
        combinatorial problem, so it is approximated by beam search or, as in CommVQ, by a \emph{learned
        encoder} that predicts the indices directly.
\end{itemize}
A useful special case is \textbf{binary additive} coding: a codebook $\mathbf{E}\in\R^{B\times m}$ and a
bit vector $\vb\in\{0,1\}^B$ decode to $\hat{\vx}=\vb\T\mathbf{E}$, i.e.\ each bit says whether to add one
codeword. The decoder is \emph{linear in the code}, and this is what makes it cheap inside attention: for
weights $a_t$,
\begin{equation}\label{eq:c13-fold}
  \sum_t a_t\,\hat{\vx}_t = \sum_t a_t\,\vb_t\T\mathbf{E} = \Bigl(\sum_t a_t\vb_t\Bigr)\T\mathbf{E}.
\end{equation}
The weighted sum is taken over bits, and the codebook is applied once, independently of the number of
tokens.

\subsection{RoPE, and why quantization must respect it}\label{sec:c13-rope}

Rotary position embedding~\citep{c01-su2024roformer} (\cref{ch:foundations}) splits a $d_k$-dimensional
query or key into $d_k/2$ pairs and rotates pair $j$ at position $t$ by the angle $t\omega_j$, with
frequencies $\omega_j=\text{base}^{-2(j-1)/d_k}$:
\[
  \bm{R}_t=\operatorname{blockdiag}\bigl(R(t\omega_1),\dots,R(t\omega_{d_k/2})\bigr),\qquad
  R(\varphi)=\begin{pmatrix}\cos\varphi&-\sin\varphi\\ \sin\varphi&\cos\varphi\end{pmatrix}.
\]
Which coordinates form a pair is a convention. The ``rotate-half'' layout of Llama-style code pairs
coordinate $j$ with $j+d_k/2$ and implements the rotation as \texttt{x * cos + rotate\_half(x) * sin} with
\texttt{rotate\_half(x) = cat(-x2, x1)}. Because $R(a)\T R(b)=R(b-a)$, the attention score depends only on
the offset:
\begin{equation}\label{eq:c13-rope-rel}
  (\bm{R}_s\vq)\T(\bm{R}_t\vk) = \vq\T\bm{R}_s\T\bm{R}_t\vk = \vq\T\bm{R}_{t-s}\vk .
\end{equation}
It helps to read each pair $(x,y)$ as a complex number $x+iy$. Then $R(\varphi)$ is multiplication by
$e^{i\varphi}$, and the matrix $aI+bJ$ with $J=\bigl(\begin{smallmatrix}0&-1\\1&0\end{smallmatrix}\bigr)$
is multiplication by $a+ib$.

\paragraph{Two bad options.} A KV quantizer can compress keys after or before the rotation. If it
stores a code for $\bm{R}_t\vk_t$ (\emph{after} RoPE), the distribution of each pair depends on the position.
High-frequency pairs are spread over all angles, and the codebook spends its codewords on rotations
instead of on the structure of the keys. If it stores a code for $\vk_t$ (\emph{before} RoPE), the codebook
sees the natural key distribution, but every cached key must be decoded and rotated at every decoding
step. That materializes full-precision keys and wastes the bandwidth the compression was meant to save.

\paragraph{The way out: codebooks that commute with rotations.} Suppose the decoder is linear,
$\hat{\vk}=D(\vec{c})$, and commutes with every $\bm{R}_t$: $\bm{R}_tD(\vec{c})=D(\bm{R}_t\vec{c})$, where the
rotation also acts on the (structured) code. Then the score is
\[
  \vq\T\bm{R}_t\hat{\vk}_t = \vq\T D(\bm{R}_t\vec{c}_t) = \bigl(D^{*}\vq\bigr)\T(\bm{R}_t\vec{c}_t),
\]
with $D^{*}$ the adjoint of $D$. The product $D^{*}\vq$ is computed once per query, and each cached key
needs only its small code rotated. Which linear maps commute with rotations? On one pair the answer is
exact.

\begin{lemma}\label{lem:c13-commute}
Let $\varphi\notin\pi\mathbb{Z}$. A real $2\times2$ matrix $X$ satisfies $XR(\varphi)=R(\varphi)X$ if and
only if $X=aI+bJ$ for some $a,b\in\R$, i.e.\ $X$ is multiplication by a complex number.
\end{lemma}
\begin{proof}
Write $R(\varphi)=\cos\varphi\,I+\sin\varphi\,J$. Every $X$ commutes with $I$, so
$XR-RX=\sin\varphi\,(XJ-JX)$, and since $\sin\varphi\ne0$ the condition is $XJ=JX$. For
$X=\bigl(\begin{smallmatrix}p&q\\r&s\end{smallmatrix}\bigr)$,
$XJ=\bigl(\begin{smallmatrix}q&-p\\s&-r\end{smallmatrix}\bigr)$ and
$JX=\bigl(\begin{smallmatrix}-r&-s\\p&q\end{smallmatrix}\bigr)$. These are equal if and only if $q=-r$ and
$s=p$, i.e.\ $X=pI+rJ$. Conversely, $aI+bJ$ commutes with $R(\varphi)$ because both are polynomials in
$J$ (in complex terms, multiplication is commutative).
\end{proof}

So a RoPE-commuting codebook is one built from $2\times2$ blocks of the form
$\bigl(\begin{smallmatrix}a&-b\\b&a\end{smallmatrix}\bigr)$, one per rotary pair. A generic real
$2\times2$ block does not commute (we confirmed this numerically, and also that the solution space of
$XR=RX$ is two-dimensional). This is the design principle of CommVQ.

\subsection{Checkpointing a recurrent state}\label{sec:c13-checkpoint}

For a recurrent layer $\mS_t=f(\mS_{t-1},\vx_t)$ with read-out $\vo_t=g(\mS_t,\vx_t)$, the state $\mS_t$
summarizes the prefix completely: all future outputs depend on $\vx_{1:t}$ only through $\mS_t$. A
\emph{checkpoint} or \emph{snapshot} is a stored copy of $\mS_b$ at some boundary $b$. It costs one
state ($H d_kd_v$ numbers per layer) and can be read later, after the live state has moved on and
partly overwritten what it held at $b$.

For ungated linear attention, $\mS_t=\sum_{s\le t}\vv_s\vk_s\T$, and snapshots taken at boundaries
$b_1<b_2<\dots$ can be organized in two ways. \emph{Prefix} (or checkpoint) snapshots are
$\mS_{b_i}$ themselves, each containing the whole prefix. \emph{Segment-local} (or restart) states
$\mS^{\text{loc}}_i=\sum_{b_{i-1}<s\le b_i}\vv_s\vk_s\T$ start from zero in each segment. They are related by
$\mS_{b_i}=\sum_{i'\le i}\mS^{\text{loc}}_{i'}$. With a decay gate the prefix form becomes
$\mS_t=(\prod_{s=b+1}^{t}\alpha_s)\mS_b+\sum_{b<s\le t}(\prod_{r=s+1}^t\alpha_r)\vv_s\vk_s\T$. The
old snapshot $\mS_b$ keeps what decay would otherwise erase. For nonlinear memories (an MLP updated by
gradient steps, \cref{ch:ttt}) the two organizations are genuinely different, because the update from
a nonzero starting state is not the update from zero plus an offset.

\paragraph{How many snapshots?} With a fixed segment length $C$ there are $\lceil T/C\rceil$ snapshots
and memory $\bigO((T/C)\,d_kd_v)$, between the $\bigO(d_kd_v)$ of an RNN and the $\bigO(T(d_k+d_v))$ of
attention. With \emph{logarithmic} segmentation the prefix is split according to the binary expansion
of its length, for example $23=16+4+2+1$. There are then at most $\lceil\log_2(T+1)\rceil$ segments, as
in a Fenwick tree~\citep{c07-fenwick1994}, and this is the structure behind log-linear attention
\citep{guo2025loglinear} (\cref{ch:structured}). The word ``checkpointing'' also names a different
technique, gradient checkpointing~\citep{c01-chen2016checkpoint}, which discards activations during
the forward pass and recomputes them for the backward pass. Here a checkpoint is part of the model's
memory at inference time.

% ===========================================================================
\section{CommVQ: commutative vector quantization of the KV cache}\label{sec:c13-commvq}
% ===========================================================================

\subsection{Motivation}

At one or two bits per number, scalar quantization has only two or four levels and the KV cache
loses too much (\cref{sec:c13-sq}). Vector quantization does much better at equal rate
(\cref{sec:c13-vq}), but two obstacles stand in the way. Encoding must be fast enough to run on every
new token, and decoding must not undo the savings. As \cref{sec:c13-rope} showed, the second obstacle
comes from RoPE: either the codebook must cover all rotations, or every key must be decoded and rotated
at every step. CommVQ~\citep{li2025commvq} removes both obstacles. Values are encoded by a small learned
network into a \emph{binary additive} code, whose linear decoder folds into the attention sum
\eqref{eq:c13-fold}. Keys are encoded with a residual codebook whose $2\times2$ blocks are complex
numbers, so by \cref{lem:c13-commute} the codebook \emph{commutes} with RoPE and attention scores can
be computed from the codes without reconstructing the keys. The survey groups CommVQ under KV-cache
compression and quantization; its one-sentence summary speaks of ``vector quantization of communication
states''~\citep{zhoubian2026memory}, but the paper's title and the repository README make clear that
``Comm'' stands for \emph{commutative} and that the object being compressed is the KV cache.

\begin{taxonomy}
The survey's Table~I does not list CommVQ; it is discussed in \S V-B (memory management and
efficiency)~\citep{zhoubian2026memory}. In terms of the three axes (our classification):
\textbf{representation} implicit (it compresses the KV cache, the working memory of attention);
\textbf{update dynamics} online for the cache contents (every new token is encoded at inference) and
offline for the codebooks and encoder (trained once); \textbf{persistence} short-term. The closest
update-rule family in Table~II is ``admission, eviction, and consolidation'', in its compressive
form: every entry is admitted but stored in compressed form.
\end{taxonomy}

\subsection{Formulation}\label{sec:c13-commvq-form}

Quantization acts per layer and per token on the concatenation of all key-value heads, a vector of
size $D=H_{kv}d_k$ ($D=1024$ for Llama-3.1-8B). Let $b\in\{1,2\}$ be the target bits per channel.

\paragraph{Values: learned encoder, binary additive codebook.} Let $\vec{\lambda}\in\R^D$ be a frozen
per-channel scale, the mean of $\abs{v}$ over calibration data. A value $\vv$ is normalized and encoded
as follows:
\begin{align}
  \tilde{\vv}&=\vv\oslash\vec{\lambda},\qquad \rho=\norm{\tilde{\vv}},\nonumber\\
  \vec{\ell} &= \mW_2\,\mathrm{GELU}\bigl(\mW_1\LayerNorm(\tilde{\vv}/\rho)+\vb_1\bigr)+\vb_2\in\R^{G\times4},
  \qquad G=Db/2, \label{eq:c13-commvq-enc}\\
  \vy_g &= \mathrm{GumbelSoftmax}_{\mathrm{hard}}(\vec{\ell}_g)\in\{\ve_0,\ve_1,\ve_2,\ve_3\},\nonumber\\
  (b_{2g},b_{2g+1}) &= (y_{g,0}+y_{g,2},\;y_{g,1}+y_{g,2}). \label{eq:c13-commvq-bits}
\end{align}
Each group makes a four-way choice that is turned into two bits: category 0 gives $(1,0)$, 1 gives
$(0,1)$, 2 gives $(1,1)$ and 3 gives $(0,0)$. A four-way softmax can express any joint distribution
over two bits, which two independent sigmoids cannot. The code $\vb\in\{0,1\}^{Db}$ is decoded with a
binary additive codebook $\mathbf{E}\in\R^{Db\times D}$, and the norm is restored:
\begin{equation}\label{eq:c13-commvq-vdec}
  \vec{x}=\vb\T\mathbf{E},\qquad \hat{\vv}=\vec{\lambda}\had\rho\,\frac{\vec{x}}{\norm{\vec{x}}}.
\end{equation}
The cost is $Db$ bits plus the scalar $\rho$ per token and layer. For $b=1$ that is one bit per
channel, a factor of 16 below fp16 (the encoder asserts \texttt{int(16 / self.quant\_bits) == compression\_ratio}).
The encoder and codebook are trained with the LLM \emph{frozen}, using the straight-through Gumbel-softmax of
\cref{sec:c13-hard} and two losses: the mean squared error between $\hat{\vv}$ and $\vv$, and the mean
squared error between each attention layer's output computed with quantized values and its original
output.

\paragraph{Keys: a residual codebook that commutes with RoPE.} Keys are quantized \emph{before}
RoPE. Each token's pre-RoPE key $\vk\in\R^D$ is divided by its norm $\rho=\norm{\vk}$ (taken over all
heads together). Within a head of dimension $d_k=128$, coordinates $j$ and $j+64$ form a rotary pair, written
as one complex number $\kappa_{h,j}=k_{h,j}+i\,k_{h,j+64}$, $j=1,\dots,P$ with $P=64$. Stage $r$ of the
residual codebook holds, for every head $h$ and pair $j$, a list of $C=64$ complex numbers
$z^{(r,h,j)}_0,\dots,z^{(r,h,j)}_{63}$. A stage code is a pair of indices $(A_r,B_r)\in\{0,\dots,63\}^2$,
i.e.\ 12 bits, \emph{shared by all 64 pairs of the head}. The decoder is
\begin{equation}\label{eq:c13-commvq-key}
  \hat{\kappa}_{h,j} = \rho\sum_{r=1}^{R}\Bigl(z^{(r,h,j)}_{A_r}+i\,z^{(r,h,j)}_{B_r}\Bigr).
\end{equation}
In real $2\times2$ form, with $M(z)=(\operatorname{Re}z)I+(\operatorname{Im}z)J$,
\[
  z_A+i\,z_B = M(z_A)\ve_1+M(z_B)\ve_2 = \begin{pmatrix}x_A-y_B\\ y_A+x_B\end{pmatrix},
\]
where $z=x+iy$. The code of a pair is the pair of unit vectors $(\ve_1,\ve_2)$ placed in the slots $A$
and $B$, and the decoder applies the blocks $M(z)$ to it. Every $M(z)$ commutes with every rotation
(\cref{lem:c13-commute}), so rotating the decoded key equals decoding a rotated code:
\begin{equation}\label{eq:c13-commvq-commute}
  R(\varphi)\bigl(M(z_A)\ve_1+M(z_B)\ve_2\bigr) = M(z_A)\bigl(R(\varphi)\ve_1\bigr)+M(z_B)\bigl(R(\varphi)\ve_2\bigr).
\end{equation}
With $R$ stages a head costs $12R$ bits for 128 channels. The training script's default is $R=21$
(252 bits, about $1.97$ bits per channel). The released 1-bit and 2-bit codebooks are downloaded
separately and fix their own number of stages. Encoding is greedy residual VQ (\cref{sec:c13-rvq}): at
each stage, pick the nearest of the $64^2=4096$ codewords and subtract it.

\paragraph{Training the key codebook by structured EM.} The codewords of one stage are
$4096$ points that are all determined by the $64$ complex numbers $z_0,\dots,z_{63}$ (per pair). They are
therefore a linear function $\vec{c}=\mT\vtheta$ of the parameters $\vtheta$ (the real and imaginary
parts of the $z_a$), with each pair of rows of $\mT$ encoding $(x_A-y_B,\,x_B+y_A)$. EM proceeds as in
\cref{sec:c13-vq}: soft assignments \eqref{eq:c13-estep} with an annealed temperature, weighted means
$\vm_c$, and then the structured M-step \eqref{eq:c13-structured-mstep}, which projects the means onto the
commuting family. For this $\mT$ one can check that $\mT\T\mT=2C\,\mI$ is diagonal (each parameter appears in
$2C$ codewords with coefficient $\pm1$, and the cross term between $x_A$ and $y_B$ from codeword $(A,B)$
cancels the one from codeword $(B,A)$). The repository asserts this.
After a stage converges, the keys are replaced by their residuals and the next stage is fitted.

\paragraph{Attention scores directly from the codes.} Let $\tilde{q}_{h,j}$ be the rotated query, in
complex form, and $\varphi_{t,j}=t\,\omega_j$ the angle of pair $j$ at key position $t$. The real dot
product of two pairs is $\operatorname{Re}(\bar a\,b)$, and RoPE multiplies the key pair by
$e^{i\varphi}$. Substituting \eqref{eq:c13-commvq-key},
\begin{equation}\label{eq:c13-commvq-score}
\begin{aligned}
  \tilde{\vq}\T\bm{R}_t\hat{\vk}_t
  &= \sum_j\operatorname{Re}\bigl(\overline{\tilde q_j}\,e^{i\varphi_{t,j}}\,\hat\kappa_{t,j}\bigr)
   = \rho_t\sum_j\operatorname{Re}\Bigl(e^{i\varphi_{t,j}}\bigl(W^{A}_{t,j}+i\,W^{B}_{t,j}\bigr)\Bigr),\\
  W^{A}_{t,j}&=\sum_r U_{r,j,A_{r,t}},\qquad U_{r,j,a}=\overline{\tilde q_j}\,z^{(r,j)}_a ,
\end{aligned}
\end{equation}
and $W^B$ likewise with the indices $B_{r,t}$. The table $U$ has $R\cdot P\cdot C$ entries and is
computed once per query, independently of $T$. Each cached key then needs $2R$ table lookups per pair,
a sum, one complex rotation and a real part. The full-precision key is never formed. Writing
$\vu=M(z)\T\tilde{\vq}$ in real coordinates, the rotation step becomes
$u^A_x\cos\varphi+u^A_y\sin\varphi+u^B_y\cos\varphi-u^B_x\sin\varphi$, which is the line that appears
in the Triton kernel (\cref{lst:c13-commvq-kernel}).

\paragraph{Value read-out from the codes.} With attention weights $a_t$ and $\rho'_t=\rho_t/\norm{\vb_t\T\mathbf{E}}$
(precomputed when the token is encoded),
\begin{equation}\label{eq:c13-commvq-vread}
  \sum_t a_t\hat{\vv}_t = \vec{\lambda}\had\Bigl(\sum_t a_t\rho'_t\,\vb_t\Bigr)\T\mathbf{E},
\end{equation}
by \eqref{eq:c13-fold}: a weighted sum over bit vectors, then one product with the codebook.

\paragraph{The recent window.} Tokens are quantized in blocks of 128. The most recent, incomplete block
is kept in full precision (\texttt{residual\_length = 128}) and quantized when it fills up. The prompt is
processed with full-precision attention, so quantization error enters only the decoding steps.

\paragraph{Cost.} Per token and layer the cache holds $Db$ value bits, $12RH_{kv}$ key bits and two
16-bit scales, against $2\cdot16D$ bits in fp16. Per decoding step and query head, the key path costs
$\bigO(RPC)$ for the table plus $\bigO(TRP)$ lookups, and the value path $\bigO(TDb)$ bit-weighted additions
plus one $Db\times d_k$ codebook product. This is not less arithmetic than fp16 attention. The gain is in
memory capacity and memory traffic, which dominate the cost of decoding.

\subsection{Algorithms}

\begin{algorithm}[t]
\caption{Training one stage-wise residual key codebook (one layer, one head), as in \texttt{quantize\_key\_cache.py}}
\label{alg:c13-commvq-em}
\begin{algorithmic}[1]
\Require pre-RoPE keys $\{\vk_n\}$ from calibration text; stages $R$; EM steps $S$; temperatures $\tau_0>\tau_1$
\State normalize: $\vk_n\gets\vk_n/\norm{\vk_n}$; reorder coordinates into rotary pairs; $s\gets1$
\For{$r=1,\dots,R$}
  \State initialize codewords from random data points; one E-step and structured M-step
  \For{$i=0,\dots,S-1$}
    \State $\tau\gets s\,\tau_0(\tau_1/\tau_0)^{i/(S-1)}$ \Comment{exponential annealing}
    \State E-step: $w_{nc}\propto\exp(-\norm{\vk_n-\vec{c}_c}/\tau)$; record the hard-assignment error
    \State keep the codebook with the lowest error; stop if the error stops improving
    \State M-step: $\vm_c\gets$ weighted means; $\vtheta\gets(\mT\T\mathbf{N}\mT)^{-1}\mT\T\mathbf{N}\vm$; $\vec{c}\gets\mT\vtheta$
  \EndFor
  \State $\vk_n\gets\vk_n-\vec{c}_{i(\vk_n)}$ for all $n$ \Comment{residual for the next stage}
  \State \textbf{if} $r\bmod 4=1$ \textbf{then} $s\gets s/10$
\EndFor
\end{algorithmic}
\end{algorithm}

\begin{algorithm}[t]
\caption{One decoding step of attention over a CommVQ cache (one query head)}
\label{alg:c13-commvq-decode}
\begin{algorithmic}[1]
\Require query $\vq$ at position $s$; codes $(A_{r,t},B_{r,t})$, scales $\rho_t,\rho'_t$, bits $\vb_t$ for $t\le T_q$; recent full-precision $(\vk_t,\vv_t)$ for $T_q<t\le s$
\State $\tilde{\vq}\gets\bm{R}_s\vq$; table $U_{r,j,a}\gets\overline{\tilde q_j}z^{(r,j)}_a$ \Comment{independent of $T$}
\For{$t=1,\dots,T_q$} \Comment{parallel over $t$}
  \State $W^A_j\gets\sum_rU_{r,j,A_{r,t}}$, $W^B_j\gets\sum_rU_{r,j,B_{r,t}}$
  \State $\text{score}_t\gets\rho_t\sum_j\operatorname{Re}(e^{it\omega_j}(W^A_j+iW^B_j))/\sqrt{d_k}$
\EndFor
\State $\text{score}_t\gets\tilde{\vq}\T\bm{R}_t\vk_t/\sqrt{d_k}$ for the recent window
\State $\va\gets\softmax(\text{score})$
\State $\vo\gets\vec{\lambda}\had(\sum_{t\le T_q}a_t\rho'_t\vb_t)\T\mathbf{E}+\sum_{t>T_q}a_t\vv_t$
\end{algorithmic}
\end{algorithm}

\subsection{Implementation walkthrough}\label{sec:c13-commvq-impl}

\begin{implnote}[title={In the code: map of \texttt{repos/UMass-Embodied-AGI\_\_CommVQ}}]
\textbf{Training pipeline} (\repofile{repos/UMass-Embodied-AGI__CommVQ/training/}): (1)
\code{collect_kv.sh} runs \code{collect_kv.py} on Llama-3.1-8B-Instruct with
\code{HuggingFaceFW/fineweb-edu} and sequence length 8192 to dump keys and values; (2)
\code{make_scale.py} computes the per-channel mean of $\abs{x}$ for every layer's keys and values and
saves \code{data/learnable_scale.pt} (the $\vec{\lambda}$ above); (3) \code{quantize_key_cache.sh} runs
\code{quantize_key_cache.py} once per layer (0 to 31) to fit the key codebooks by EM; (4)
\code{finetune/llama3.1_8b_int1.sh} (and siblings for 2 bits, Llama-3.1-70B, Mistral-7B and
Qwen-7B) runs \code{finetune.py} with \texttt{--is\_stage1 True --use\_mse\_loss True --use\_vq\_loss True}.
In stage 1 every parameter whose name does not contain \code{value_vq_model} is frozen, and
\code{learnable_scale} is frozen as well.
\textbf{Models} (\repofile{repos/UMass-Embodied-AGI__CommVQ/commvq/}): \code{modeling_llama_training.py}
(training, with the two MSE losses), \code{modeling_llama_evaluation.py} (simulated quantization:
codes are decoded back to dense keys and values before attention), \code{modeling_llama_triton.py}
(attention computed from the codes with the Triton kernels in \code{triton_kernels.py}). The
compressors are in \code{compress_training.py}, \code{compress_evaluation.py} and \code{compress_triton.py}.
\end{implnote}

The heart of the key codebook is the matrix $\mT$ and the structured M-step
\eqref{eq:c13-structured-mstep}. The constants at the top of the script are \texttt{KV\_DIM = 1024},
\texttt{N\_BITS = 12}, \texttt{M\_GROUP = 128} (one head), \texttt{RESIDUAL\_NUM = 21}, \texttt{N\_STEPS = 100}, and
temperatures from $10^{-3}$ to $10^{-5}$.

\begin{lstlisting}[caption={Excerpt from \texttt{repos/UMass-Embodied-AGI\_\_CommVQ/training/quantize\_key\_cache.py}},label={lst:c13-commvq-mstep}]
T = torch.zeros(2*POW_2_N_BITS_HALF*POW_2_N_BITS_HALF, 2*POW_2_N_BITS_HALF).cuda()
for A in range(POW_2_N_BITS_HALF):
    for B in range(POW_2_N_BITS_HALF):
        T[2*(A*POW_2_N_BITS_HALF+B), 2*A] = 1
        T[2*(A*POW_2_N_BITS_HALF+B), 2*B+1] = -1
        T[2*(A*POW_2_N_BITS_HALF+B)+1, 2*B] = 1
        T[2*(A*POW_2_N_BITS_HALF+B)+1, 2*A+1] = 1
# ...
def M_step(tensors, clustering_centers, soft_assignments):
    assert len(clustering_centers) == CLUSTERING_CENTER_NUM
    # m: new mean of each cluster
    # N_tensors: number of tensors in each cluster
    N_tensors = torch.sum(soft_assignments, dim=0).unsqueeze(-1).repeat(1, 2) + 1e-8
    weighted_sums = torch.matmul(soft_assignments.T, tensors)
    m = weighted_sums / N_tensors[:, :1]
    N_tensors = N_tensors.flatten().unsqueeze(0)
    mat = (T.T * N_tensors) @ T
    A_mat = (torch.linalg.inv(mat) @ T.T) * N_tensors
    B_mat = m.view(m.shape[0], M_GROUP // 2, 2).permute(0, 2, 1).flatten(0, 1)
    theta = A_mat @ B_mat
    clustering_centers = (T @ theta).view(CLUSTERING_CENTER_NUM, 2, M_GROUP // 2).permute(0, 2, 1).flatten(1)
    return clustering_centers, theta
\end{lstlisting}

\begin{implnote}[title={In the code: reading the M-step}]
Row pair $2(64A+B),\,2(64A+B)+1$ of \code{T} produces the codeword $(x_A-y_B,\;x_B+y_A)$ from
$\vtheta=(x_0,y_0,x_1,y_1,\dots,x_{63},y_{63})$, i.e.\ $z_A+i\,z_B$. \code{mat} is
$\mT\T\mathbf{N}\mT$ and \code{A_mat} is $(\mT\T\mathbf{N}\mT)^{-1}\mT\T\mathbf{N}$, so \code{theta} is
\eqref{eq:c13-structured-mstep}. \code{B_mat} rearranges the weighted means so that each of the 64 rotary
pairs of the head becomes one column. All pairs share $\mT$ and the soft counts and differ only in their
data, so one matrix product solves 64 least-squares problems. Before the loop, the data are reshaped
with \texttt{view(N, 8, 2, 64).transpose(2, 3)}, which interleaves coordinates $j$ and $j+64$ of each head
(the rotate-half pairs). The E-step (\code{E_step}) uses \code{torch.cdist}, the \emph{unsquared}
distance, inside \texttt{softmax(-dist / temperature)}. As $\tau\to0$ this still selects the nearest
codeword. \code{handle_single} runs the stages of \cref{alg:c13-commvq-em}.
\end{implnote}

The value encoder's quantizer and decoder are short:

\begin{lstlisting}[caption={Excerpt from \texttt{repos/UMass-Embodied-AGI\_\_CommVQ/commvq/compress\_evaluation.py} (\texttt{ValueCompressor})},label={lst:c13-commvq-value}]
    def quantize(self, x):
        B, S, E = x.shape
        x = x.reshape(B, S, self.n_e, self.code_range)
        x = torch.nn.functional.gumbel_softmax(x, tau=1.0, hard=True)
        x[..., :2] += x[..., 2:3]
        x = x[..., :2].flatten(2)
        return x

    def decode(self, x, prescale):
        x = torch.matmul(x, self.embedding)
        postscale = torch.norm(x, dim=-1, keepdim=True)
        x = x / (postscale + self.eps)
        x = x * (prescale + self.eps)
        x = x * (self.learnable_scale + self.eps)
        return x
\end{lstlisting}

Here \texttt{code\_range = 4}, \texttt{n\_e = feat\_dim * quant\_bits // 2} groups, and \code{embedding} is the
codebook $\mathbf{E}$ with \texttt{n\_e * 2} rows. \code{pack} stores the bits eight to a \code{uint8}.
The key path stores \code{uint16} codes of shape \texttt{[R, H\_kv, tokens]}, and \code{decode2} recovers the
two 6-bit halves as \texttt{x // 64} and \texttt{x \% 64}. In the Triton path the score
\eqref{eq:c13-commvq-score} is computed in two kernels. \code{calculate_QC_T_kernel} builds the table
$\vu=M(z)\T\tilde{\vq}$ for every query head, pair, stage and codeword. The kernel below then gathers,
sums over stages, rotates and reduces:

\begin{lstlisting}[caption={Excerpt from \texttt{repos/UMass-Embodied-AGI\_\_CommVQ/commvq/triton\_kernels.py}},label={lst:c13-commvq-kernel}]
    for r_i in tl.static_range(0, R):
        c_start = hx * R * N + r_i * N + token_offset
        c = tl.load(key_code_ptr + c_start)  # FIXME: do I need mask?
        c1x = c // 64
        c2x = c % 64
        c1_offset = codebook_start[None, :] + h_i * R * C * 2 + r_i * C * 2 + c1x[:, None] * 2
        c2_offset = codebook_start[None, :] + h_i * R * C * 2 + r_i * C * 2 + c2x[:, None] * 2
        c1x = tl.load(QC_T_ptr + c1_offset)
        c1y = tl.load(QC_T_ptr + c1_offset + 1)
        c2x = tl.load(QC_T_ptr + c2_offset)
        c2y = tl.load(QC_T_ptr + c2_offset + 1)
        acc_1X += c1x
        acc_1Y += c1y
        acc_2X += c2x
        acc_2Y += c2y
    sin = tl.load(sin_ptr + sin_start)
    cos = tl.load(cos_ptr + cos_start)
    acc = tl.sum(acc_1X * cos + acc_1Y * sin + acc_2Y * cos - acc_2X * sin, axis=1)
    output_offset = (hx * H_RATIO + h_i) * N + token_offset
    acc = acc * prescale
    tl.store(output_ptr + output_offset, acc, mask=mask)
\end{lstlisting}

\begin{implnote}[title={In the code: the decoding attention}]
In \code{LlamaFlashAttention2.forward} of \code{modeling_llama_triton.py}, the first call (prefill) runs
FlashAttention on full-precision tensors and stores the codes of the first
$\lfloor T/128\rfloor\cdot128$ tokens. The key norm is stored already divided by $\sqrt{d_k}$ (\texttt{pack(..., factor=math.sqrt(self.head\_dim))}),
so the kernels need no separate scaling, and the value prescale is stored already divided by the postscale
($\rho'_t$ above, via \code{cache_postscale}). At each later step the query is rotated, the two kernels
above produce the scores of the quantized tokens, a dense product handles the recent full-precision
window, and \code{attn_weights_mul_value_triton} computes \eqref{eq:c13-commvq-vread}: it unpacks the
value bits inside the kernel, accumulates $\sum_ta_t\rho'_t\vb_t$ with \code{tl.dot}, and multiplies by
the codebook (and $\vec{\lambda}$) once. That kernel asserts a 1024-bit value code (``Currently only
1-bit quantization is supported to use Triton kernels''), matching the README's note that the
memory-saving kernels support Llama-3.1-8B at 1 bit. The query-head axis is handled by repeating the
8 key-value heads' codebooks 4 times (grouped-query attention with 32 query heads).
One caution for readers who run this file: in the mirrored snapshot, the branch that compresses the
recent window once it reaches 128 tokens still contains debugging statements (\texttt{print("???")} and
\texttt{import pdb; pdb.set\_trace()}, under a \texttt{FIXME: [DEBUG ONLY]} comment), so long generations
with the Triton model stop there. The evaluation model, \code{modeling_llama_evaluation.py}, has no such
statements.
\end{implnote}

\subsection{Reference implementation}

\Cref{lst:c13-commvq-ref} implements the two decoders, greedy residual encoding of a key, and the
score computed from codes. For simplicity it normalizes each head's key separately and uses one head.

\begin{lstlisting}[caption={Reference implementation (ours): additive VQ of a KV vector with a RoPE-commuting key codebook},label={lst:c13-commvq-ref}]
import torch

def encode_key(k, Z):
    """k: [P, 2] pre-RoPE key as rotary pairs; Z: [R, P, C] complex codebooks.
    Greedy residual VQ; stage codeword (A, B) is z_A + i z_B (12 bits if C = 64)."""
    rho = k.norm()
    x = torch.complex(k[:, 0], k[:, 1]) / rho
    R, P, C = Z.shape
    codes = []
    for r in range(R):
        cw = Z[r].T[:, None, :] + 1j * Z[r].T[None, :, :]   # [C, C, P]
        A, B = divmod(int((x - cw).abs().pow(2).sum(-1).argmin()), C)
        codes.append((A, B))
        x = x - cw[A, B]                                     # residual
    return codes, rho

def decode_key(codes, Z, rho):                               # -> complex [P]
    return rho * sum(Z[r, :, A] + 1j * Z[r, :, B] for r, (A, B) in enumerate(codes))

def score_from_codes(q_rot, codes, Z, rho, phi):
    """q_rot: complex [P] rotated query; phi: [P] angles t * omega_j of the key.
    Equals Re(conj(q_rot) * e^{i phi} * decode_key(...)).sum() without decoding."""
    U = torch.conj(q_rot)[None, :, None] * Z                 # table, [R, P, C]
    WA = sum(U[r, :, A] for r, (A, B) in enumerate(codes))
    WB = sum(U[r, :, B] for r, (A, B) in enumerate(codes))
    return rho * torch.real(torch.exp(1j * phi) * (WA + 1j * WB)).sum()

def decode_value(bits, E, rho, lam, eps=1e-6):               # bits: [n_bits] in {0,1}
    x = bits @ E                                             # sum of selected codewords
    return lam * rho * x / (x.norm() + eps)

def read_values(a, bits, E, rho_post, lam):
    """sum_t a_t * decode_value(bits_t) with rho_post = rho / ||bits @ E||, in one pass."""
    return lam * (((a * rho_post) @ bits) @ E)
\end{lstlisting}

\paragraph{Numerical checks.} In NumPy we verified: (i) \cref{lem:c13-commute} (the solution space of
$XR=RX$ is two-dimensional, of the form $aI+bJ$); (ii) $R(\theta)\,\mathrm{decode}(c)=\mathrm{decode}(R(\theta)c)$
for \eqref{eq:c13-commvq-commute} (error $4\times10^{-16}$), and that the decoded codeword is
$(x_A-y_B,\,x_B+y_A)$, the row pattern of \code{T}; (iii) scores from the table, in the kernel's real form
and in the complex form \eqref{eq:c13-commvq-score}, match decode-rotate-dot with three stages (error
$4\times10^{-15}$); (iv) the value read-out \eqref{eq:c13-commvq-vread}; (v) the 4-way-to-2-bit map;
(vi) $\mT\T\mT$ is diagonal and the closed-form M-step minimizes the weighted distortion (no random
perturbation of $\vtheta$ improved it).

\begin{results}
\textbf{Source:} \repofile{repos/UMass-Embodied-AGI__CommVQ/README.md}. The README gives no numbers.
It states that CommVQ ``achieves strong performance across a wide range of benchmarks while
significantly reducing memory overhead'', releases Llama-3.1-8B checkpoints (value and key codebooks)
for 1-bit and 2-bit CommVQ, provides evaluation scripts for LongBench, InfiniteBench, needle-in-a-haystack
and memory measurement, and notes of the Triton kernels for real memory savings: ``Currently
supports LLaMA-3.1 8B with 1-bit quantization''. The survey describes the method as ``significantly
reducing KV storage [\dots] with minimal performance degradation''~\citep{zhoubian2026memory}. See the paper
for numbers.
\end{results}

\subsection{Discussion}

\paragraph{Strengths.} CommVQ reaches one or two bits per channel by spending its bits on vectors
rather than numbers, and it keeps the decode cost low by matching the codebook's algebra to RoPE's. The
LLM itself is not retrained: the encoder and codebooks are fitted to a frozen model, so the method
applies to existing checkpoints.

\paragraph{Limitations.} (i) Codebooks are fitted per model and per layer on calibration data, and a
domain shift at inference is not corrected. (ii) Encoding costs a small MLP per value and a 4096-way
nearest-neighbour search per key stage (on blocks of 128 tokens). (iii) The arithmetic of attention does
not shrink, only the memory traffic, so the gain depends on how memory-bound decoding is. (iv) In the
released code the fast kernels cover one model and one bit width.

\paragraph{Relations.} CommVQ is orthogonal to eviction~\citep{c13-zhang2023h2o} and admission (HAM):
those decide \emph{which} entries to keep, and CommVQ decides \emph{how} to store them, so the two could
be combined. Low-rank latent caches~\citep{deepseek2024v2} also compress the cache, but along the
feature dimension and by changing the architecture. The idea of making a compressed representation
commute with the model's structure, so that computation can stay in the compressed domain, applies
beyond RoPE.

% ===========================================================================
\section{Memory Caching: recurrent memory that grows}\label{sec:c13-mc}
% ===========================================================================

\subsection{Motivation}

A recurrent memory has a fixed size, so it must overwrite or decay old content to make room for new
content (\cref{ch:linear,ch:delta}). Attention never overwrites, but it pays for every token.
Memory Caching~\citep{behrouz2026memorycaching} sits between the two. It runs an ordinary recurrent
memory, but at the end of every segment of the sequence it saves a copy of the memory state
(\cref{sec:c13-checkpoint}), and later tokens read from all saved copies as well as from the live
state. The survey summarizes it as storing ``segment-level snapshots of recurrent hidden memory'' and
aggregating them later, ``creating a middle ground between fixed-size recurrent states and fully
token-level Transformer caches''~\citep{zhoubian2026memory}. The memory grows with the number of
segments, not with the number of tokens.

\begin{pitfall}[title={The code is unofficial}]
There is no official code. The mirrored repository \repo{kmccleary3301__memory_caching} is a community
implementation. Its README says plainly: ``This is not official author code.'' Its documentation
separates a ``mechanism-faithful'' implementation, which it claims, from ``paper parity'', which it marks
as ``Blocked'' (\repofile{repos/kmccleary3301__memory_caching/docs/reproduction_report.md}). The README's
BibTeX entry for the paper also gives an author list (``Chandra, ...'') that does not match this book's
bibliography, which follows the survey. Everything below describes the mechanism \emph{as implemented
there}, using the paper's names for the aggregation rules. The original paper may differ in details.
\end{pitfall}

\begin{taxonomy}
\textbf{Representation:} implicit (recurrent states and their snapshots). \textbf{Update dynamics:}
online. \textbf{Persistence:} long-term (survey Table~I: ``Cached Recurrent State Snapshots'').
\textbf{Update rule:} Table~II does not list Memory Caching. In our reading it combines a state-transition
update (the backbone recurrence) with a retention rule: every segment's final state is kept and none
is evicted~\citep{zhoubian2026memory}.
\end{taxonomy}

\subsection{Formulation}\label{sec:c13-mc-form}

\paragraph{Ingredients.} A \emph{memory backend} has a state $\mM$, an update
$\mM_t=U(\mM_{t-1};\vk_t,\vv_t)$ and a read-out $A(\mM,\vq)$. For linear attention,
$U(\mM;\vk,\vv)=\mM+\vv\vk\T$ and $A(\mM,\vq)=\mM\vq$. The sequence is cut into segments
$\mathcal{S}_1,\mathcal{S}_2,\dots$, either of constant length $C$ or of lengths given by the binary
expansion of $T$, largest first ($23\to16,4,2,1$). Let $\sigma(t)$ be the segment containing $t$.

\paragraph{Online state and snapshots.} Inside segment $i$ the layer updates an online state token by
token. It starts either from the previous segment's final state (\code{checkpoint}, the default) or
from the initial state (\code{restart}). It also accumulates a segment context, the running sum of the
segment's keys:
\begin{equation}\label{eq:c13-mc-online}
  \mM_t = U(\mM_{t-1};\vk_t,\vv_t),\qquad \vec{c}^{\,\mathrm{on}}_t=\sum_{s\in\mathcal{S}_{\sigma(t)},\,s\le t}\vk_s .
\end{equation}
At the end of segment $i$ the layer caches the snapshot $\mM^{(i)}$ (the final online state) and the
context $\vec{c}^{(i)}$ (the sum of all keys in the segment).

\paragraph{Aggregation.} Token $t$ reads from the $n=\sigma(t)-1$ cached snapshots and the online
state. A separate projection $\vu_t=\mW_U\vx_t$ (or $\vq_t$ itself) scores the segments by their
contexts. The four rules implemented are:
\begin{align}
  \text{Residual (RM):}\quad &\vy_t = A(\mM_t,\vq_t)+\sum_{i=1}^{n}A(\mM^{(i)},\vq_t),\label{eq:c13-mc-rm}\\
  \text{Gated residual (GRM):}\quad &\vy_t = \gamma^{\mathrm{on}}_tA(\mM_t,\vq_t)+\sum_{i=1}^{n}\gamma^{(i)}_tA(\mM^{(i)},\vq_t),\label{eq:c13-mc-grm}\\
  &(\gamma^{(1)}_t,\dots,\gamma^{(n)}_t,\gamma^{\mathrm{on}}_t)=\softmax\bigl(\vu_t\T\vec{c}^{(1)},\dots,\vu_t\T\vec{c}^{(n)},\vu_t\T\vec{c}^{\,\mathrm{on}}_t\bigr),\nonumber\\
  \text{Memory soup:}\quad &\vy_t = A\Bigl(\gamma^{\mathrm{on}}_t\mM_t+\textstyle\sum_{i}\gamma^{(i)}_t\mM^{(i)},\;\vq_t\Bigr),\label{eq:c13-mc-soup}\\
  \text{Sparse selective (SSC):}\quad &\text{GRM over the online state and the top-}k\text{ snapshots by }\vu_t\T\vec{c}^{(i)}.\label{eq:c13-mc-ssc}
\end{align}
GRM is a softmax ``attention over segments'': the keys are segment summaries, the query is $\vu_t$, and
the values are the snapshots' read-outs. Soup mixes the \emph{states} before reading. SSC bounds the
read cost by visiting only $k$ snapshots (default $k=2$).

\begin{proposition}\label{prop:c13-soup}
If the read-out is linear in the state, $A(\mM,\vq)=\mM\vq$, then Memory Soup equals GRM.
\end{proposition}
\begin{proof}
$A(\sum_i\gamma_i\mM^{(i)},\vq)=(\sum_i\gamma_i\mM^{(i)})\vq=\sum_i\gamma_i\mM^{(i)}\vq=\sum_i\gamma_iA(\mM^{(i)},\vq)$.
\end{proof}
For a deep memory, such as a two-layer MLP updated by gradient steps, $A$ is nonlinear in the
parameters, and mixing parameters differs from mixing outputs. Only then is Soup a separate method.

\begin{proposition}\label{prop:c13-rm-restart}
For ungated linear attention with \code{restart} initialization, RM equals plain linear attention.
\end{proposition}
\begin{proof}
With restart, $\mM^{(i)}=\sum_{s\in\mathcal{S}_i}\vv_s\vk_s\T$ and the online state is
$\sum_{s\in\mathcal{S}_{\sigma(t)},s\le t}\vv_s\vk_s\T$. Their sum is $\sum_{s\le t}\vv_s\vk_s\T=\mS_t$, and the
read-out is linear, so \eqref{eq:c13-mc-rm} gives $\mS_t\vq_t$.
\end{proof}
So caching pays off only when something breaks this additivity: decay or a delta rule (the snapshot
keeps what the live state has since overwritten), a nonlinear memory, or input-dependent gating. With
\code{checkpoint} initialization and RM, the same linear backend gives
$\vy_t=\sum_{s\le t}(1+m_s)\,\vv_s\vk_s\T\vq_t$, where $m_s$ is the number of cached prefix snapshots
that contain $s$. Older tokens are counted more often, which is one reason the gated rules exist.

\paragraph{Backends.} Four memories are implemented: \code{linear} (above); \code{swla} with $c=2$,
$\mM_t=\alpha\mM_{t-1}+\beta\vv_{t-1}\vk_{t-1}\T+\lambda\vv_t\vk_t\T$; \code{dla}, a per-head two-layer
MLP (GELU) whose parameters take one gradient step per token on $-\inner{\mathrm{MLP}(\vk_t)}{\vv_t}$
(the \code{dot} objective) or $\norm{\mathrm{MLP}(\vk_t)-\vv_t}^2$, optionally with momentum; and
\code{titans}, a deep memory with a momentum buffer in the style of \cref{ch:ttt}. For a one-layer
linear ``MLP'' the \code{dot} objective has gradient $-\vv_t\vk_t\T$, so one step of size $\eta$ is
$\mM_t=\mM_{t-1}+\eta\vv_t\vk_t\T$, linear attention again. DLA is its deep generalization.

\paragraph{Cost.} With constant segments of length $C$ the layer holds $\lceil t/C\rceil$ states,
i.e.\ $\bigO((t/C)\,d_kd_v)$ numbers per head. RM, GRM and Soup read every snapshot, $\bigO((t/C)\,d_kd_v)$
work per token and $\bigO(T^2d_kd_v/C)$ for a sequence: still quadratic, but with a $1/C$ factor and with no
softmax over tokens. Logarithmic segments give $\bigO(\log T)$ states. SSC scores all contexts
($\bigO((t/C)d_k)$, cheap) but reads only $k$ states.

\subsection{Algorithm}

\begin{algorithm}[t]
\caption{Memory Caching layer, per-token form (as in the unofficial implementation)}
\label{alg:c13-mc}
\begin{algorithmic}[1]
\Require projections $\vq_t,\vk_t,\vv_t,\vu_t$; backend $(U,A)$; segments; aggregation rule
\State cache $\gets[\,]$; $\mM\gets\mM_0$
\For{each segment $\mathcal{S}_i$}
  \If{restart}
    \State $\mM\gets\mM_0$
  \EndIf
  \State $\vec{c}\gets\bm{0}$
  \For{$t\in\mathcal{S}_i$}
    \State $\mM\gets U(\mM;\vk_t,\vv_t)$; $\vec{c}\gets\vec{c}+\vk_t$
    \State $\vy_t\gets$ aggregate $A(\cdot,\vq_t)$ over cache $\cup\{(\mM,\vec{c})\}$ by RM, GRM, Soup or SSC
  \EndFor
  \State append $(\operatorname{copy}(\mM),\operatorname{copy}(\vec{c}))$ to cache \Comment{snapshot}
\EndFor
\State \Return $\mW_O[\vy_1;\dots;\vy_T]$
\end{algorithmic}
\end{algorithm}

The same loop serves decoding: the layer's state is the online memory, the online context and the list
of snapshots. A training implementation for real models would use the chunkwise algorithms of
\cref{ch:linear} inside each segment. The unofficial code loops over tokens in Python, which is meant for
clarity, not speed.

\subsection{Implementation walkthrough}\label{sec:c13-mc-impl}

\begin{implnote}[title={In the code: map of \texttt{repos/kmccleary3301\_\_memory\_caching}}]
\begin{itemize}
  \item \repofile{repos/kmccleary3301__memory_caching/src/memory_caching/layer.py}:
        \code{MemoryCachingLayer} with projections \code{q_proj}, \code{k_proj}, \code{v_proj},
        \code{u_proj}, \code{o_proj} ($d\to d$, split into heads), \code{_forward_impl} (the loop of
        \cref{alg:c13-mc}), and \code{_residual_aggregate}, \code{_grm_or_soup_aggregate},
        \code{_ssc_aggregate}.
  \item \code{config.py}: \code{MCConfig} defaults \code{backend="linear"}, \code{aggregation="grm"},
        \code{segmentation="constant"}, \code{segment_size=256}, \code{state_init_mode="checkpoint"},
        \code{ssc_top_k=2}, \code{use_q_as_u=False}, \code{softmax_temperature=1.0},
        \code{detach_cached_states=False}.
  \item \code{contracts.py}: \texttt{SegmentCache(mem\_state, seg\_context, seg\_len)} and the
        \code{MemoryBackend} protocol (\code{init_state}, \code{update}, \code{apply}); a
        \code{MixableMemoryBackend} adds \code{mix_states} for Soup.
  \item \code{segmentation.py} (\code{constant_segments}, \code{logarithmic_segments}),
        \code{backends/} (\code{linear.py}, \code{swla.py}, \code{dla.py}, \code{titans.py}),
        \code{baselines/loglinear_pp.py} (\code{LogLinearPP}, which fixes \code{aggregation="grm"}
        and \code{segmentation="logarithmic"}; compare \cref{ch:structured}), and
        \code{tests/test_paper_equations.py} (analytic checks of RM, GRM, Soup and SSC).
\end{itemize}
Shapes: per-token tensors are \texttt{[B, H, d\_h]}; the linear state is \texttt{[B, H, d\_h, d\_h]}.
Snapshots are taken with \code{tree_clone} (or \code{tree_detach_clone} when
\code{detach_cached_states=True}), so by default gradients flow from later segments back into earlier
snapshots, i.e.\ backpropagation through time across segments.
\end{implnote}

\begin{lstlisting}[caption={Excerpt from \texttt{repos/kmccleary3301\_\_memory\_caching/src/memory\_caching/layer.py}},label={lst:c13-mc-grm}]
    def _grm_or_soup_aggregate(
        # ...
    ) -> Tensor:
        contexts = [c.seg_context for c in cached] + [online_context]
        scores = self._context_scores(u_t, contexts)
        weights = torch.softmax(scores, dim=-1)

        if self.config.aggregation == "soup" and self._backend_supports_state_mixing:
            states = [c.mem_state for c in cached] + [online_state]
            mixed_state = self.backend.mix_states(states, weights)
            return self.backend.apply(mixed_state, q_t)
        if self.config.aggregation == "soup" and not self.config.allow_output_mixture_fallback:
            raise ValueError(
                "aggregation='soup' requires a backend that supports state mixing; "
                "set allow_output_mixture_fallback=True to enable output-space fallback."
            )

        responses = [self.backend.apply(c.mem_state, q_t) for c in cached]
        responses.append(self.backend.apply(online_state, q_t))
        stacked_responses = torch.stack(responses, dim=2)  # [B,H,S,Dh]
        return torch.einsum("bhs,bhsd->bhd", weights, stacked_responses)
\end{lstlisting}

The elided keyword arguments are \code{q_t}, \code{u_t}, \code{online_state}, \code{cached} (a list of
\code{SegmentCache}) and \code{online_context}. \code{_context_scores} stacks the contexts to
\texttt{[B, H, S, d\_h]}, contracts them with \code{u_t} by an \code{einsum}, and divides by the temperature.
In \code{_forward_impl} the online context is reset to zero at each segment start and incremented by
\code{k_t} after each update, and at the segment end a \code{SegmentCache} with the state, the context and
the segment length is appended. SSC
masks all but the top-$k$ cached scores with $-\infty$, appends the online score, takes one softmax, and
skips the read-out of unselected snapshots.

\begin{pitfall}[title={Sign conventions in the \texttt{titans} backend}]
\code{backends/titans.py} offers two conventions. The default, \code{update_convention="paper"}, computes
$\mathbf{s}_t=\beta\mathbf{s}_{t-1}-\eta\nabla\ell$ and $\mM_t=\alpha\mM_{t-1}-\mathbf{s}_t$. With zero momentum this is
$\mM_t=\alpha\mM_{t-1}+\eta\nabla\ell$, a step that \emph{increases} the inner loss. The repository's own test
(\repofile{repos/kmccleary3301__memory_caching/tests/test_titans_update_convention.py}) hand-derives exactly
this: a target of $1$, a prediction of $0$, and a bias that moves to a negative value. The alternative
\code{"gradient_descent"} descends. Titans as defined in \cref{ch:ttt} descends the inner loss.
\code{docs/TITANS_NOTES.md} says the \code{paper} mode mirrors ``the written recurrence form''. Check the
sign before you use this backend. The \code{dla} backend always descends.
\end{pitfall}

\subsection{Reference implementation}

\begin{lstlisting}[caption={Reference implementation (ours): segment-snapshot memory caching with gated aggregation (linear backend, GRM)},label={lst:c13-mc-ref}]
import torch

def memory_caching(q, k, v, u, seg=256, init="checkpoint", temp=1.0):
    """q, k, u: [T, d_k]; v: [T, d_v]. Linear memory M <- M + v k^T,
    snapshots at segment ends, gated residual (softmax) aggregation."""
    T, d_k = k.shape
    snaps, ctxs, out = [], [], []              # cached states and contexts
    M = k.new_zeros(v.shape[1], d_k)
    for start in range(0, T, seg):
        if init == "restart":
            M = k.new_zeros(v.shape[1], d_k)
        c = k.new_zeros(d_k)                    # running segment summary
        for t in range(start, min(start + seg, T)):
            M = M + torch.outer(v[t], k[t])     # online write
            c = c + k[t]
            scores = torch.stack([u[t] @ ci for ci in ctxs + [c]]) / temp
            gamma = torch.softmax(scores, dim=0)              # [n_seg + 1]
            reads = torch.stack([Mi @ q[t] for Mi in snaps + [M]])
            out.append(gamma @ reads)                         # GRM read-out
        snaps.append(M.clone()); ctxs.append(c.clone())       # snapshot
    return torch.stack(out)
\end{lstlisting}

We checked in NumPy, with a re-implementation that follows \code{layer.py} (read, not imported):
(i) Soup equals GRM for the linear backend (\cref{prop:c13-soup}, difference $4\times10^{-15}$); (ii) RM with
\code{restart} equals plain linear attention (\cref{prop:c13-rm-restart}); (iii) SSC with $k$ at least the
number of segments equals GRM; (iv) \code{checkpoint} snapshots equal the prefix states; (v) the number of
snapshots is $\lceil T/C\rceil$; (vi) the logarithmic segmentation of $23$ is $(16,4,2,1)$ and of $1000$ is
$(512,256,128,64,32,8)$.

\begin{results}
\textbf{Source:} \repofile{repos/kmccleary3301__memory_caching/README.md} and
\repofile{repos/kmccleary3301__memory_caching/docs/}. The unofficial repository reports no paper-scale
results. Its status table lists ``Full table-level paper parity: Blocked by missing baselines''.
\code{docs/TRAINING_PARITY_TABLE.md} records that its own runs used sequence length 128 for 4 steps,
against target profiles of 4096 to 16384 tokens and 1000 to 10000 steps. The survey states qualitatively
that Memory Caching creates ``a middle ground between fixed-size recurrent states and fully token-level
Transformer caches''~\citep{zhoubian2026memory}. See the paper for numbers.
\end{results}

\subsection{Discussion}

\paragraph{Strengths.} Memory Caching gives any recurrent memory a knob, the segment length, that
interpolates between an RNN ($C=T$) and something close to attention ($C=1$, \cref{exr:c13-mc}). It is
backend-agnostic, and its gated rules are themselves a small attention over segments, so retrieval
across segments is content-based.

\paragraph{Limitations.} Snapshots are large: one $d_k\times d_v$ state per head per segment, against
$d_k+d_v$ numbers per token for a KV cache. Caching beats a KV cache in memory only if $C>d_kd_v/(d_k+d_v)$,
i.e.\ $C>32$ for $d_k=d_v=64$. Reads cost a state product per snapshot. The segment context is a plain sum
of keys, a coarse summary when segments are long.

\paragraph{Relations.} With logarithmic segments and GRM, Memory Caching resembles log-linear
attention (\cref{ch:structured}), which keeps $\bigO(\log T)$ states per query. The repository makes the
link explicit through its \code{LogLinearPP} baseline. Compared with HAM, which keeps \emph{tokens}
selected by surprise, Memory Caching keeps \emph{states} selected by position. CommVQ could compress the
snapshots as well.

% ===========================================================================
\section{Bottlenecked Transformers: consolidating the KV cache}\label{sec:c13-bottleneck}
% ===========================================================================

\subsection{Motivation}

Every method so far keeps cache entries as written, or drops or compresses them. In a decoder-only
Transformer, the entry for position $s$ is a causal function of $x_{1:s}$, computed once and never
revisited. It cannot discard details that later prove irrelevant, and it cannot absorb what later
tokens reveal about it. Bottlenecked Transformers~\citep{oomerjee2026bottlenecked} add a small
auxiliary Transformer, the \emph{Cache Processor}, that periodically rewrites KV entries \emph{in place}
at reasoning-step boundaries. The motivation comes from neuroscience. \emph{Consolidation} stabilizes
newly formed memory traces~\citep{c13-mcgaugh2000}, and \emph{reconsolidation} makes a recalled
memory labile again so that it can be updated~\citep{c13-nader2000}. It also comes from information
theory: the \emph{information bottleneck} principle~\citep{c13-tishby2000ib} holds that a good
representation $Z$ of an input $X$ for predicting $Y$ solves
\begin{equation}\label{eq:c13-ib}
  \min_{p(z\mid x)}\; I(X;Z)-\beta\,I(Z;Y),
\end{equation}
i.e.\ it keeps as little about the input as possible while staying predictive of the target. Here
$X$ is the context read so far, $Z$ the KV cache, and $Y$ the tokens still to be generated. A cache that is
written once and never revised cannot move towards this trade-off; a processor that rewrites it can. The
survey describes the method as ``periodic processing of cached representations, using an auxiliary
bottleneck module to consolidate or rewrite parts of the working memory at structured generation
boundaries''~\citep{zhoubian2026memory}.

\begin{taxonomy}
The survey's Table~I does not list Bottlenecked Transformers. Table~II places them under
``admission, eviction, and consolidation'' (with ``irreversible loss and retrieval bias'' as main risks).
In \S V-B the survey adds that, ``from the taxonomy perspective, this remains closer to short-term
implicit working memory than to long-term explicit storage'', because the memory being changed is
the model's own cache~\citep{zhoubian2026memory}. On the three axes: \textbf{representation} implicit,
\textbf{update dynamics} online (rewrites during generation), \textbf{persistence} short-term.
\end{taxonomy}

\subsection{Formulation}\label{sec:c13-bt-form}

No code is public. From public descriptions: the backbone LLM is frozen; the Cache Processor consists of
small Transformer blocks, one aligned to each backbone layer; it is invoked when a reasoning step ends,
detected by a newline token; it performs non-causal, in-place rewrites of the entries of the most recent
step (consolidation) and of a small set of earlier entries selected by top-$k$ attention mass
(reconsolidation); and it preserves the length and dimension of the cache. \emph{One way to formalize
this is the following; the paper's exact parameterization may differ.}

Let $\mK^{(\ell)},\mV^{(\ell)}$ be layer $\ell$'s cache and $b_1<b_2<\cdots$ the step boundaries. At
boundary $b_j$ define the recent window and the recalled set
\begin{equation}\label{eq:c13-bt-sets}
  \mathcal{W}_j=\{b_{j-1}+1,\dots,b_j\},\qquad
  \mathcal{P}^{(\ell)}_j=\operatorname{TopK}_{s\le b_{j-1}}\Bigl(\sum_{t\in\mathcal{W}_j}\sum_{h}a^{(\ell,h)}_{t,s}\Bigr),
\end{equation}
where $a^{(\ell,h)}_{t,s}$ is the attention weight of query $t$ on key $s$ in head $h$. These weights
were already computed while the step was generated, so they can be accumulated during decoding. The
processor of layer $\ell$ reads the selected entries as a set of tokens, $\vx_s=[\vk^{(\ell)}_s;\vv^{(\ell)}_s]$
for $s\in\mathcal{W}_j\cup\mathcal{P}^{(\ell)}_j$, applies non-causal self-attention and an MLP among
them, and writes back a residual update:
\begin{equation}\label{eq:c13-bt-rewrite}
  \bigl[\vk^{(\ell)}_s;\vv^{(\ell)}_s\bigr]\;\leftarrow\;\bigl[\vk^{(\ell)}_s;\vv^{(\ell)}_s\bigr]
  + \mathrm{Proc}^{(\ell)}\bigl(\{\vx_{s'}\}_{s'\in\mathcal{W}_j\cup\mathcal{P}^{(\ell)}_j}\bigr)_s,
  \qquad s\in\mathcal{W}_j\cup\mathcal{P}^{(\ell)}_j .
\end{equation}
With a zero-initialized output projection, the rewrite starts as the identity and the model starts as
the backbone. The same trick appears in AMOR. All later tokens attend to the rewritten entries. Because
the backbone is frozen, the processor must produce entries that the backbone's own queries can still
read. This is a strong constraint and keeps the rewrites ``in distribution''. The training objective and
how positions are handled for keys that already carry RoPE are details for which we refer to the paper.

\begin{algorithm}[t]
\caption{Generation with periodic KV consolidation (our formalization)}\label{alg:c13-bt}
\begin{algorithmic}[1]
\Require frozen backbone; processors $\mathrm{Proc}^{(1..L)}$; budget $k$
\State prefill the prompt; $b\gets$ prompt length
\Repeat
  \State generate the next token $x_t$ with the backbone; accumulate attention mass on positions $\le b$
  \If{$x_t$ is a newline}
    \For{$\ell=1,\dots,L$} \Comment{parallel over layers}
      \State $\mathcal{W}\gets\{b+1,\dots,t\}$; $\mathcal{P}\gets$ top-$k$ earlier positions by mass
      \State rewrite $(\mK^{(\ell)},\mV^{(\ell)})$ on $\mathcal{W}\cup\mathcal{P}$ by \eqref{eq:c13-bt-rewrite}
    \EndFor
    \State $b\gets t$; reset the accumulated mass
  \EndIf
\Until{end of sequence}
\end{algorithmic}
\end{algorithm}

\paragraph{Cost.} The cache keeps its size, so the memory cost is unchanged. Per boundary and layer, the
processor attends among $m+k$ entries ($m=\abs{\mathcal{W}_j}$), costing $\bigO((m+k)^2d)$. Amortized over
the $m$ tokens of the step, this is $\bigO((m+k)^2d/m)$ per token, small next to attention over a long
cache when steps are short and $k$ is small. The processor's parameters are an additional, small, fixed cost.

\begin{results}
\textbf{Source:} no public code was found (\repofile{README.md}, section~4). The survey makes the
qualitative point that ``attention memory can be managed not only by dropping, quantizing, or paging
entries, but also by periodically transforming them into a more compact or stable representation
before subsequent reasoning steps''~\citep{zhoubian2026memory}. The authors report, on mathematical
reasoning benchmarks with several backbone LLMs, consistent gains over vanilla Transformers and over
baselines augmented with pause tokens~\citep{c13-goyal2024pause}. See the paper for numbers.
\end{results}

\subsection{Discussion}

\paragraph{Strengths.} Consolidation adds latent computation without adding tokens. A chain of thought
spends compute by writing text, pause tokens~\citep{c13-goyal2024pause} by writing placeholder tokens,
and the Cache Processor by revising the memory itself. The rewrite is bounded in scope (the last step
plus $k$ recalled entries), so the overhead is predictable. The backbone stays frozen.

\paragraph{Limitations.} (i) In-place rewriting breaks a property that serving systems rely on: the KV
entries of a prefix are no longer a function of the prefix alone. Prefix caching and block sharing in
paged KV memory~\citep{kwon2023paged} (\cref{ch:foundations}) need copy-on-write for rewritten blocks.
(ii) Boundaries are defined by newlines, a heuristic tied to step-by-step reasoning formats. (iii) As
with every consolidation scheme, a rewrite can destroy detail that would have been needed later; the
survey's ``irreversible loss'' risk applies.

\paragraph{Relations.} Like CommVQ, the processor changes the \emph{content} of cache entries rather
than their number. CommVQ does so to save bits, with a fixed, offline-trained map; the Cache Processor
does so to improve prediction, with a learned map that depends on context. Like HAM and AMOR it uses
attention statistics or other signals to decide \emph{which} entries to touch.

% ===========================================================================
\section{Comparison and discussion}\label{sec:c13-compare}
% ===========================================================================

\Cref{tab:c13-compare} lines up the five systems on the questions of this chapter: which signal
drives the decision, what is admitted, compressed or rewritten, how memory grows with the sequence length
$T$, and what the mechanism costs.

\begin{table}[t]
\centering\small
\renewcommand{\arraystretch}{1.25}
\begin{tabularx}{\textwidth}{@{}>{\raggedright\arraybackslash}p{2.1cm}>{\raggedright\arraybackslash}X>{\raggedright\arraybackslash}X>{\raggedright\arraybackslash}p{2.6cm}>{\raggedright\arraybackslash}X@{}}
\toprule
System & Signal & Admitted / compressed / rewritten & Memory growth & Overhead \\
\midrule
AMOR & normalized output entropy vs.\ EMA threshold (or fitted router) & nothing is dropped; attention \emph{runs} only for uncertain tokens & $\bigO(1)$ backbone $+$ $\bigO(KTd)$ for $K$ attention blocks & LM-head product per block and token (router: $\approx$\,1/250 of it) \\
HAM (our formalization) & RNN prediction error vs.\ threshold (or learned router) & tokens the RNN predicts poorly are admitted to the KV cache & $\bigO(d_kd_v)+\bigO(rT(d_k+d_v))$, $r$ = admission rate & one extra state-vector product per token \\
CommVQ & none (all tokens) & every key and value compressed to codes ($\approx$1--2 bits/channel) & $\bigO(T)$ with a $\approx$8--16$\times$ smaller constant & encoder MLP and codeword search at write; code-domain attention at read \\
Memory Caching & segment boundaries; router $\vu_t\T\vec{c}^{(i)}$ over segments & recurrent state snapshots at segment ends & $\bigO((T/C)d_kd_v)$, or $\bigO(\log T)$ states & one state read per snapshot (SSC: $k$) \\
Bottlenecked Transformers & newline boundaries; top-$k$ attention mass & recent-step and recalled entries rewritten in place & $\bigO(T)$, unchanged & processor over $m+k$ entries per boundary \\
\bottomrule
\end{tabularx}
\caption{The five systems of this chapter. ``Our formalization'' marks rows whose mechanism details
come from \cref{sec:c13-ham-form} rather than released code.}
\label{tab:c13-compare}
\end{table}

\paragraph{Where to intervene.} The five systems act at different points in the life of a memory
entry. AMOR acts at \emph{read} time (does this token query the high-resolution memory?). HAM acts at
\emph{write} time (does this token get a verbatim entry?). CommVQ changes the \emph{representation} (how
many bits an entry takes). Memory Caching changes the \emph{granularity} (whole states instead of
tokens). Bottlenecked Transformers \emph{revise} entries after the fact. These choices are largely
independent. A HAM-style layer could quantize its admitted entries with a CommVQ codebook, snapshot its
recurrent state at segment ends, and consolidate the admitted entries at step boundaries.

\paragraph{Signals.} Three kinds of signal appear. \emph{Model uncertainty} (AMOR's output entropy)
says how unsure the prediction is. \emph{Memory error} (HAM's innovation) says what the compact memory
fails to hold, and is the more direct measure of what a second memory should store. \emph{Usage} (the
attention mass of Bottlenecked Transformers, and of heavy-hitter eviction~\citep{c13-zhang2023h2o}) says
what has been useful so far. All three are computed from quantities the model produces anyway, but each
has a known failure mode: irreducible uncertainty, surprising noise, and delayed relevance respectively.

\paragraph{Compute versus memory.} The systems save different resources. AMOR saves attention compute
but, in its released decoder, not KV memory. HAM and Memory Caching save memory by storing less.
CommVQ saves memory and memory bandwidth but not arithmetic. Bottlenecked Transformers spend extra
compute for quality at fixed memory. When reading claims about ``efficient memory'', always ask which
resource is meant.

\paragraph{Open problems.} Calibrating thresholds across layers, heads and input distributions;
recovering from wrong admission decisions (a token not admitted is gone); keeping ragged,
data-dependent caches efficient on hardware built for dense tensors; and avoiding a mismatch between
training (dense masked computation, as in AMOR and the parallel form of HAM) and inference (sparse
computation). \Cref{ch:synthesis} returns to these questions across the whole book.

\begin{summarybox}
\begin{itemize}
  \item Fixed-ratio hybrids decide once where attention goes. \textbf{Adaptive routing} decides per token
        which memory a token uses, driven by a signal: output entropy (AMOR) or the recurrent memory's
        prediction error (HAM). Threshold policies control the rate, and an EMA of quantiles keeps it stable.
  \item \textbf{AMOR} appends entropy-gated attention blocks to a recurrent language model. The gate is
        detached and parameter-free, dense-masked training equals sparse inference, the released decoder
        still stores every key and value, and a fitted router replaces the costly LM-head entropy.
  \item \textbf{HAM} lets an RNN see every token and admits only poorly predicted tokens to a KV cache.
        The cache grows with novelty, deduplicates what the RNN already knows, and is controlled by one
        threshold.
  \item \textbf{Managing working memory} uses several independent levers: fewer heads, fewer tokens,
        fewer bits, states instead of tokens, paging, and rewriting.
  \item \textbf{CommVQ} compresses keys and values to 1--2 bits per channel with vector quantization: a
        learned encoder with a binary additive codebook for values, and a residual codebook of complex
        $2\times2$ blocks for keys. These blocks are exactly the matrices that commute with RoPE, so scores
        and value sums are computed from the codes. The key codebook is trained by EM with a structured
        least-squares M-step.
  \item \textbf{Memory Caching} snapshots a recurrent memory at segment ends and aggregates snapshots
        with residual, gated, soup or sparse-selective rules. For linear memories Soup equals GRM, and
        residual-with-restart collapses to linear attention. The only code is unofficial.
  \item \textbf{Bottlenecked Transformers} rewrite recent and recalled KV entries in place at reasoning-step
        boundaries with a small processor on a frozen backbone, motivated by memory consolidation and
        the information bottleneck.
\end{itemize}
\end{summarybox}

\section*{Exercises}

\begin{exercise}[Entropy and fire rates]\label{exr:c13-entropy}
(a) Derive $\mathrm{H}(\softmax(\vz))=\operatorname{lse}(\vz)-\sum_ip_iz_i$ and show that adding a constant to
all logits leaves it unchanged. (b) Compute $\bar{\mathrm{H}}$ for $\vec{p}=(0.5,0.25,0.25,0)$ over
$\abs{\mathcal{V}}=4$. (c) If the normalized entropies are $\mathcal{N}(\mu,\sigma^2)$, show that AMOR's
adaptive threshold $\tau=\mu+\kappa\sigma$ fires on a fraction $1-\Phi(\kappa)$ of positions, and find the
$\kappa$ that gives a 10\% fire rate. (d) Explain why a skewed entropy distribution breaks this
calculation, and propose an EMA estimator of the $(1-r)$-quantile that targets a fire rate $r$ directly.
\end{exercise}

\begin{exercise}[Why AMOR caches every key]\label{exr:c13-amor}
(a) Prove \cref{prop:c13-amor-sparse} for multi-head attention with RoPE. (b) Suppose the decoder stored
keys and values only at firing positions. Show by a two-token example that a firing query's output then
differs from the dense-masked training computation. (c) With fire rate $f$, $K$ blocks, width $d$ and
vocabulary $\abs{\mathcal{V}}$, write the per-token decoding cost of the native gate and of attention at
context length $t$, and find the $t$ above which the attention saving exceeds the gate's cost.
\end{exercise}

\begin{exercise}[Surprise-gated admission with decay]\label{exr:c13-ham}
Consider the delta-rule memory \eqref{eq:c13-ham-rnn} with $\beta_t=1$, constant decay $\alpha<1$ and
orthonormal keys. (a) Show that an item written at time $t$ is recalled at time $t+n$ as $\alpha^n\vv$ if
no other write uses its key, so its normalized error is $(1-\alpha^n)^2$. (b) With threshold $\tau$, after how
many steps is a repeated item admitted again? (c) Implement \cref{lst:c13-ham-ref} in NumPy, reproduce
\cref{ex:c13-ham-dedup}, and plot cache size against $\tau$ for a stream with 60\% repeats.
(d) Construct a sequence where a token is not admitted but is needed later, and discuss a fix based on
delayed admission (keeping a short buffer of candidates).
\end{exercise}

\begin{exercise}[Codebooks that commute with RoPE]\label{exr:c13-commute}
(a) Let $\bm{R}_t$ be block diagonal with blocks $R(t\omega_j)$, $j=1,\dots,P$, and distinct
frequencies $\omega_j\in(0,\pi)$. Show that a $2P\times2P$ matrix commutes with $\bm{R}_t$ for all integers $t$ if and
only if it is block diagonal with blocks $a_jI+b_jJ$. (Hint: for the blocks $X_{jk}$, consider
$R(t\omega_j)X_{jk}=X_{jk}R(t\omega_k)$; use \cref{lem:c13-commute} for $j=k$.) (b) Derive
\eqref{eq:c13-commvq-score} from \eqref{eq:c13-commvq-key}, and check it numerically against
decode-rotate-dot. (c) Count the bits per channel of CommVQ keys with $R$ stages and of values with
$b$ bits, including the two 16-bit scales, for $D=1024$ and $H_{kv}=8$.
\end{exercise}

\begin{exercise}[Structured EM]\label{exr:c13-em}
(a) Derive \eqref{eq:c13-structured-mstep} and show it reduces to weighted means when $\mT=\mI$. (b) For the
CommVQ matrix $\mT$ with $C$ complex parameters, show that $\mT\T\mT=2C\,\mI$, and explain why
$\mT\T\mathbf{N}\mT$ is not diagonal for general soft counts $\mathbf{N}$. (c) Implement annealed EM with the
structured M-step for $C=8$ on random unit vectors in $\R^4$ (two rotary pairs) and compare its distortion
with unconstrained $k$-means using the same number of codewords. How much does the RoPE constraint cost?
\end{exercise}

\begin{exercise}[Memory Caching in closed form]\label{exr:c13-mc}
For the linear backend: (a) prove \cref{prop:c13-soup,prop:c13-rm-restart}. (b) Take segment length
$C=1$, \code{restart} initialization and GRM. Show that
$\vy_t=\sum_{s\le t}\gamma_{t,s}(\vk_s\T\vq_t)\,\vv_s$ with $\gamma_{t,\cdot}=\softmax_s(\vu_t\T\vk_s)$. Compare
this with softmax attention: what plays the role of the attention logits, and what extra factor scales each
value? (c) Show that snapshots use less memory than a KV cache for the same tokens if and only if
$C>d_kd_v/(d_k+d_v)$. (d) Run \cref{lst:c13-mc-ref} with \code{checkpoint} and \code{restart} on a
recall task of your choice and explain the difference.
\end{exercise}
